\documentclass[11pt,a4paper]{article}
\usepackage[utf8]{inputenc}
\usepackage[T1]{fontenc}
\usepackage[spanish,es-noshorthands]{babel}
\usepackage{amsmath,amssymb}
\usepackage{booktabs,array}
\usepackage{tikz}
\usepackage[margin=2.5cm]{geometry}

\usepackage{xcolor}
\usepackage[colorlinks=true,linkcolor=blue!50!black]{hyperref}

\newenvironment{idea}{\par\begin{center}\begin{lrbox}{\ideabox}\begin{minipage}{0.9\linewidth}}{\end{minipage}\end{lrbox}\fcolorbox{blue!40!black}{blue!5}{\usebox{\ideabox}}\end{center}}
\newsavebox{\ideabox}

\title{Sección 2.3: Divergencias de Bregman\\Notas de estudio}
\author{Plan de tesis, Capítulo 2}
\date{}

\begin{document}
\maketitle
\tableofcontents

\part{Sección 2.3.1: Definición y propiedades}

\section{Qué es una divergencia}

Al agrupar datos, cada grupo se resume con un \emph{representante}, y hace falta
medir cuánto se pierde al reemplazar cada dato $x$ por su representante $y$. Esa
medida de pérdida es una divergencia.

Se parece a una distancia, pero no lo es: en general no es simétrica y no cumple
la desigualdad triangular. Por eso se la llama divergencia.

La distancia euclídea al cuadrado, $(x-y)^2$, trata igual un error de 10 viajes
en una celda de 100 y en una de 10\,000. Las divergencias de Bregman permiten
construir medidas que sí distinguen esos casos: se elige una función $\phi$ según
el tipo de dato y de ella sale la medida.

\section{El caso de un solo número}

\subsection{Qué es $\phi$ y qué es $\phi'$}

$\phi$ es una función \emph{estrictamente convexa}: una curva en forma de U, sin
tramos rectos. $\phi'$ es su \emph{derivada}: para cada punto $y$ dice qué tan
inclinada está la curva en ese punto, es decir, su \emph{pendiente}.

Ejemplo con $\phi(x)=x^2$, cuya derivada es $\phi'(x)=2x$:
\begin{center}
\begin{tabular}{cccc}
\toprule
Punto $y$ & $\phi(y)$ (altura de la curva) & $\phi'(y)$ (pendiente) & Lectura \\
\midrule
$0$ & $0$ & $0$ & la curva está plana \\
$1$ & $1$ & $2$ & sube 2 por cada unidad \\
$3$ & $9$ & $6$ & más empinada \\
\bottomrule
\end{tabular}
\end{center}

\subsection{La fórmula se construye en tres pasos}

\begin{enumerate}
  \item \textbf{La tangente en $y$.} Es la recta que toca la curva solo en $y$ y
  tiene su misma inclinación. Parte de la altura $\phi(y)$ y sube $\phi'(y)$ por
  cada unidad que avanzas. Al moverte de $y$ a $x$ (distancia $x-y$), su altura
  en $x$ es
  \[
    \text{tangente}(x)=\phi(y)+\phi'(y)\,(x-y).
  \]
  \item \textbf{La curva en $x$} tiene altura $\phi(x)$.
  \item \textbf{La divergencia} es cuánto queda la curva por encima de la
  tangente en $x$:
  \[
    d_\phi(x,y)=\phi(x)-\bigl[\phi(y)+\phi'(y)(x-y)\bigr].
  \]
\end{enumerate}
Quitando el corchete se obtiene la fórmula de la tesis:
\begin{equation}\label{eq:b1}
  d_\phi(x,y)=\phi(x)-\phi(y)-\phi'(y)\,(x-y).
\end{equation}
\begin{idea}
No hay nada más detrás de la fórmula: es una resta entre la curva y la tangente.
\end{idea}

\subsection{Ejemplo numérico}

Con $\phi(x)=x^2$, representante $y=1$ y dato $x=3$:
\begin{itemize}
  \item Curva en $x$: $\phi(3)=9$.
  \item Tangente en $x$: $\phi(1)+\phi'(1)(3-1)=1+2\cdot 2=5$.
  \item Divergencia: $9-5=\mathbf{4}$, que es $(3-1)^2$. Es la línea gruesa
  vertical de la figura de la tesis.
\end{itemize}
En general, con $\phi(x)=x^2$:
$x^2-y^2-2y(x-y)=x^2-2xy+y^2=(x-y)^2$.

\begin{center}
\begin{tikzpicture}[xscale=1.6, yscale=0.3]
  \draw[->] (-0.3,0) -- (3.9,0);
  \draw[->] (0,-1.5) -- (0,11.5);
  \draw[thick, domain=-0.25:3.3, samples=60] plot (\x,{\x*\x})
    node[above] {\small $\phi(x)=x^2$};
  \draw[dashed, domain=-0.25:3.3] plot (\x,{2*\x-1})
    node[below right] {\small tangente en $y$};
  \draw[dotted] (1,0) -- (1,1);
  \draw[dotted] (3,0) -- (3,5);
  \draw[line width=1.6pt, red!70!black] (3,5) -- (3,9);
  \node[right] at (3,7) {\small $d_\phi(x,y)=4$};
  \node[circle, fill, inner sep=1.3pt] at (1,1) {};
  \node[circle, fill, inner sep=1.3pt] at (3,9) {};
  \node[circle, fill, inner sep=1.3pt] at (3,5) {};
  \node[left] at (3,9) {\small 9};
  \node[left] at (3,5) {\small 5};
  \node[below] at (1,0) {\small $y=1$};
  \node[below] at (3,0) {\small $x=3$};
\end{tikzpicture}
\end{center}

\subsection{Por qué nunca es negativa}

Porque $\phi$ es una U: siempre queda \emph{por encima} de cualquiera de sus
tangentes. Por eso curva $-$ tangente $\ge 0$, y vale $0$ solo cuando $x=y$
(ahí curva y tangente se tocan). Convexidad implica divergencia no negativa.

\section{El caso de vectores y matrices}

\subsection{Una matriz es un vector largo}

Una matriz de conteos $2\times 2$ (origen $\times$ hora), por ejemplo
$\begin{pmatrix}100&20\\30&50\end{pmatrix}$, se aplana en un vector con 4
componentes: $x=(100,20,30,50)$. Lo mismo se hace con el representante $y$ (la
aproximación por bloques). Cada celda es una componente.

\subsection{Qué cambia}

\begin{itemize}
  \item \textbf{$\phi$ recibe un vector} y devuelve un número. Ya no es una curva
  sino una superficie (un \guillemotleft{}cuenco\guillemotright{} en más dimensiones).
  \item \textbf{Hay una pendiente por componente.} En una superficie puedes
  moverte en cada dirección y la inclinación es distinta en cada una. El vector
  de pendientes es el \emph{gradiente}
  \[
    \nabla\phi(y)=\Bigl(\tfrac{\partial\phi}{\partial y_1},\dots,
    \tfrac{\partial\phi}{\partial y_p}\Bigr),
  \]
  donde cada entrada es la derivada de $\phi$ respecto a esa componente dejando
  las demás fijas. Es la generalización de $\phi'$.
  \item \textbf{La tangente es un plano} (hiperplano). Su altura en $x$ es
  $\phi(y)+\sum_k (x_k-y_k)\,\partial\phi/\partial y_k$.
  \item La suma $\sum_k a_k b_k$ se escribe $\langle a,b\rangle$ y se llama
  producto interno: multiplica componente a componente y suma.
\end{itemize}
La idea es la misma, curva $-$ plano tangente:
\begin{equation}\label{eq:bn}
  d_\phi(x,y)=\phi(x)-\phi(y)-\langle x-y,\nabla\phi(y)\rangle .
\end{equation}
Con una sola componente, \eqref{eq:bn} coincide con \eqref{eq:b1}.

\subsection{Ejemplo numérico con la matriz}

Sea $\phi(x)=\sum_k x_k^2$, de modo que $\partial\phi/\partial y_k=2y_k$. Con
$y=(60,60,60,60)$ y $x=(100,20,30,50)$:
\begin{itemize}
  \item $\phi(x)=10\,000+400+900+2\,500=13\,800$.
  \item $\phi(y)=4\cdot 3\,600=14\,400$.
  \item $\nabla\phi(y)=(120,120,120,120)$.
  \item $x-y=(40,-40,-30,-10)$, cuya suma es $-40$.
  \item $\langle x-y,\nabla\phi(y)\rangle=120\cdot(-40)=-4\,800$.
  \item $d=13\,800-14\,400-(-4\,800)=\mathbf{4\,200}$.
\end{itemize}
Comprobación: $\sum_k (x_k-y_k)^2=1\,600+1\,600+900+100=4\,200$. Coincide con la
distancia euclídea al cuadrado.

\subsection{Por qué se puede sumar celda por celda}

Si $\phi$ es una suma de una misma función $\psi$ aplicada a cada celda,
$\phi(x)=\sum_k\psi(x_k)$, entonces $\partial\phi/\partial y_k=\psi'(y_k)$ solo
depende de la celda $k$. Al sustituir, todo se separa:
\[
  d_\phi(x,y)=\sum_k\Bigl[\psi(x_k)-\psi(y_k)-\psi'(y_k)(x_k-y_k)\Bigr]
  =\sum_k d_\psi(x_k,y_k).
\]
Por eso la discrepancia entre dos matrices se calcula como la suma de la
divergencia escalar de cada celda.

\section{Propiedades}

\begin{center}
\renewcommand{\arraystretch}{1.25}
\begin{tabular}{@{}>{\raggedright\arraybackslash}p{5.2cm}>{\raggedright\arraybackslash}p{9.4cm}@{}}
\toprule
\textbf{Propiedad} & \textbf{Por qué} \\
\midrule
$d\ge 0$, y $d=0$ solo si $x=y$ & La tangente está bajo la curva y $\phi$ es estrictamente convexa. \\
\midrule
Convexa en el primer argumento $x$ & $\phi(x)$ es convexa y el resto es lineal en $x$. En el segundo argumento no lo es en general. \\
\midrule
No simétrica: $d(x,y)\ne d(y,x)$ & La tangente se traza en $y$ (el representante), no en $x$ (el dato). \\
\midrule
No cumple la desigualdad triangular & Por eso no es una métrica y se llama divergencia. \\
\midrule
Aditiva por componentes & Si $\phi(x)=\sum\psi(x_k)$, entonces $d(x,y)=\sum d(x_k,y_k)$. Permite sumar celda por celda. \\
\bottomrule
\end{tabular}
\end{center}

\section{Los tres casos de la tabla y por qué esos $\phi$}

No hay nada mágico en esos tres: \textbf{$\phi$ es una elección de modelado}. Se
elige la curva según cómo se comporta el error en los datos. Cada $\phi$
corresponde a un supuesto distinto sobre cómo varía el error con la magnitud del
dato.

\begin{center}
\renewcommand{\arraystretch}{1.3}
\footnotesize
\begin{tabular}{@{}>{\raggedright\arraybackslash}p{2.6cm}>{\raggedright\arraybackslash}p{3.6cm}>{\raggedright\arraybackslash}p{5.6cm}>{\raggedright\arraybackslash}p{2.6cm}@{}}
\toprule
$\phi(x)$ & $d_\phi(x,y)$ & Qué supone sobre el error & Tipo de dato \\
\midrule
$x^2$ & $(x-y)^2$ & Constante en todas partes (varianza constante, como una gaussiana) & Medidas continuas \\
\midrule
$x\log x-x$ & $x\log\frac{x}{y}-(x-y)$ & Crece con el tamaño de la celda (varianza proporcional a la media, como Poisson) & Conteos \\
\midrule
$\sum_k x_k\log x_k$ & $\sum_k x_k\log\frac{x_k}{y_k}$ & Igual que el anterior, pero con datos que suman 1 & Distribuciones de probabilidad \\
\bottomrule
\end{tabular}
\end{center}

\subsection{Cómo sale la I-divergencia}

Con $\phi(x)=x\log x-x$, la derivada es $\phi'(x)=\log x$ (la derivada de
$x\log x$ es $\log x+1$ y la de $-x$ es $-1$, que se cancelan). Sustituyendo en
\eqref{eq:b1}:
\begin{align*}
  d(x,y)&=x\log x-x-(y\log y-y)-\log y\,(x-y)\\
        &=x\log x-x-y\log y+y-x\log y+y\log y\\
        &=x\log\frac{x}{y}-(x-y).
\end{align*}
El caso de Kullback--Leibler es el mismo cálculo sobre el símplex: como
$\sum_k x_k=\sum_k y_k=1$, el término $-(x-y)$ se anula al sumar.

\subsection{Por qué la I-divergencia pondera por magnitud}

Una diferencia de 10 viajes (dato 110, representante 100):
\[
  d=110\ln 1.1-10=10.484-10\approx\mathbf{0.48}.
\]
La misma diferencia en una celda grande (dato 10\,010, representante 10\,000):
\[
  d=10\,010\ln 1.001-10\approx\mathbf{0.005}.
\]
La distancia euclídea da 100 en ambos casos. Atención: la divergencia no es
simétrica; con dato 90 y representante 100 el primer valor sería $0.52$, por eso
conviene decir siempre cuál es el dato y cuál el representante.

La razón de fondo es que, cerca de $y$, la divergencia se comporta como
\[
  d_\phi(x,y)\approx\tfrac12\,\phi''(y)\,(x-y)^2 .
\]
La segunda derivada de $\phi$ es el \guillemotleft{}peso\guillemotright{} que se da al error en cada punto:
para $x^2$ es constante ($\phi''=2$); para $x\log x-x$ es $\phi''=1/x$, que
disminuye cuando $x$ crece. Ese es el comportamiento de un modelo de Poisson,
donde la varianza crece con la media.

\subsection{Por qué estos tres en la tesis}

\begin{itemize}
  \item \textbf{Euclídea:} referencia clásica ($k$-means), sirve para comparar.
  \item \textbf{I-divergencia:} la que se usa de verdad, porque los datos son
  conteos de viajes; corresponde exactamente al modelo de Poisson
  (Sección~2.3.3).
  \item \textbf{Kullback--Leibler:} caso particular de la I-divergencia con datos
  que suman 1; es la que aparece en la literatura de co-agrupamiento.
\end{itemize}
Existen más divergencias de Bregman (Itakura--Saito, Mahalanobis, \dots), pero la
tesis se queda con las tres que necesita.

\section{Homogeneidad de la I-divergencia}

Si dato y representante se multiplican por una misma constante $c>0$, la
divergencia queda multiplicada por $c$:
\[
  d_I(cx,cy)=cx\log\frac{x}{y}-c(x-y)=c\,d_I(x,y).
\]
Esto se usa más adelante para comparar conteos y promedios (sección \guillemotleft{}Escala de
los datos\guillemotright{}).

\section{De dónde viene Kullback--Leibler}

KL (Kullback--Leibler) es una divergencia pensada para \textbf{comparar dos
distribuciones de probabilidad}, es decir, dos listas de números que suman 1.
Responde a esta pregunta: si la realidad es $x$ y se la describe con el modelo
$y$, ¿cuánto se pierde?
\[
  KL(x\,\|\,y)=\sum_k x_k\log\frac{x_k}{y_k}.
\]

\subsection{Cómo sale de Bregman}

Se elige $\phi(x)=\sum_k x_k\log x_k$. Su derivada respecto a cada componente es
$\log x_k+1$. Sustituyendo en la definición:
\[
  d(x,y)=\sum_k x_k\log x_k-\sum_k y_k\log y_k-\sum_k (x_k-y_k)(\log y_k+1).
\]
Se desarrolla el último término, $\sum_k (x_k-y_k)\log y_k+\sum_k (x_k-y_k)$. Al
restarlo, los términos con $y\log y$ se cancelan y queda
\[
  d(x,y)=\sum_k x_k\log\frac{x_k}{y_k}-\sum_k (x_k-y_k).
\]
Si $x$ e $y$ son distribuciones, ambas suman 1, así que
$\sum_k (x_k-y_k)=1-1=0$ y el último término desaparece. Queda exactamente KL.
Por eso KL es la I-divergencia cuando los datos suman 1: es el mismo cálculo, solo
que aquí el término $-(x-y)$ se anula.

\subsection{Qué mide, con un ejemplo}

La realidad $x$ es un dado cargado y el modelo $y$ es un dado justo:
\begin{center}
\begin{tabular}{cccc}
\toprule
Cara & 1 & 2 & 3 y 4 \\
\midrule
Real $x$ & 0.5 & 0.5 & 0 \\
Modelo $y$ & 0.25 & 0.25 & 0.25 \\
\bottomrule
\end{tabular}
\end{center}
\[
  KL(x\|y)=0.5\ln\frac{0.5}{0.25}+0.5\ln\frac{0.5}{0.25}=\ln 2\approx 0.693.
\]
Las caras 3 y 4 no aportan porque $x=0$ ahí y se usa la convención
$0\log 0=0$. Si $y=x$, todos los $\log(x/y)$ valen 0 y KL $=0$; cuanto más se
parece el modelo a la realidad, menor es KL. Otra lectura equivalente: KL es el
número extra de bits que se gastan al codificar los datos con el código del modelo
$y$ en lugar del código óptimo para $x$.

\subsection{Por qué no es simétrica}

Con los mismos números, $KL(y\|x)$ falla porque $x$ vale 0 donde $y$ no: aparece
$\log(0.25/0)$, que es infinito. Entonces $KL(x\|y)\ne KL(y\|x)$, igual que ocurre
con las divergencias de Bregman en general.

\begin{idea}
KL compara dos distribuciones. Sale de Bregman con $\phi=\sum x\log x$. Vale 0 si
son iguales y crece cuanto más se alejan.
\end{idea}

\section{Cuándo se usa cada divergencia}

Cada una cumple un papel distinto en la tesis. Solo una es la que se aplica a los
datos; las otras dos están por referencia y por teoría.

\begin{center}
\renewcommand{\arraystretch}{1.3}
\begin{tabular}{@{}>{\raggedright\arraybackslash}p{3.0cm}>{\raggedright\arraybackslash}p{6.2cm}>{\raggedright\arraybackslash}p{5.0cm}@{}}
\toprule
\textbf{Divergencia} & \textbf{Papel en la tesis} & \textbf{Dónde aparece} \\
\midrule
\textbf{I-divergencia} & Es la que se aplica a los datos reales. & Base de co-agrupamiento (C5 adoptada, C2 como variante) y vínculo con Poisson \\
\midrule
\textbf{Kullback--Leibler} & Conecta con la literatura del co-agrupamiento por información mutua. & Función objetivo declarada por Liu et al. (2021), la pérdida de información mutua \\
\midrule
\textbf{Euclídea} & Es el caso clásico y la referencia para comparar. & K-means, varianza, bloques constantes (Cho et al., 2004) \\
\bottomrule
\end{tabular}
\end{center}

\subsection{I-divergencia: la que se usa con los datos}

La matriz de co-ocurrencia son \textbf{conteos de viajes}, no medidas continuas.
Para conteos, la I-divergencia es la apropiada por dos razones: pondera el error
según el tamaño de la celda (una diferencia de 10 viajes pesa mucho en una celda
de 100 y casi nada en una de 10\,000) y es exactamente la divergencia del modelo
de Poisson, de modo que minimizarla equivale a maximizar la verosimilitud de
Poisson (Parte~\ref{sec:parte-expfam}).

\subsection{KL: el puente con la literatura}

KL es la I-divergencia cuando los datos suman 1. Se necesita porque la función
objetivo que declaran Liu et al.\ es la pérdida de información mutua
$I(R;T)-I(\hat R;\hat T)$, y la información mutua se define con KL. Sin KL no se
podría explicar la diferencia entre lo que declara ese trabajo (base C5) y lo que
parece implementar su código (base C2).

\subsection{Euclídea: la referencia}

No es la apropiada para conteos, pero cumple tres papeles: es el punto de partida
didáctico (la que todo el mundo conoce), es la divergencia con la que se explican
primero K-means, la media como mejor representante y la varianza, y es la que usa
el co-agrupamiento clásico de bloques constantes.

\begin{idea}
I-divergencia: los datos. KL: la literatura de información mutua. Euclídea: la
intuición y la comparación.
\end{idea}

% =====================================================================
\part{Sección 2.3.2: La media como mejor representante}

\section{Qué pregunta responde}

Cada grupo de datos se resume con \textbf{un solo número}, el representante. La
pregunta es: ¿qué número elegir para que, en promedio, el error al reemplazar cada
dato por él sea el menor posible? La respuesta de esta sección es que, para
\emph{cualquier} divergencia de Bregman, el mejor representante es la media del
grupo.

Esto importa porque hace posible un algoritmo tipo K-means: se asigna cada dato al
representante más cercano y luego se recalcula cada representante como la media de
su grupo. Sin esta propiedad no habría una forma simple de recalcular los
centroides.

\section{La ecuación, símbolo por símbolo}

\[
  \mathrm{E}[d_\phi(X,s)]=\mathrm{E}[d_\phi(X,\mu)]+d_\phi(\mu,s).
\]

\subsection{El escenario}

Un grupo de datos: las celdas $1$, $2$ y $6$ de un bloque. Se quiere resumirlo con
un solo número.

\begin{center}
\renewcommand{\arraystretch}{1.3}
\begin{tabular}{@{}>{\raggedright\arraybackslash}p{2.6cm}>{\raggedright\arraybackslash}p{6.4cm}>{\raggedright\arraybackslash}p{5.2cm}@{}}
\toprule
\textbf{Símbolo} & \textbf{Qué significa} & \textbf{En el ejemplo} \\
\midrule
$X$ & Un dato del grupo, tomado al azar & Puede ser 1, 2 o 6, cada uno con probabilidad $1/3$ \\
\midrule
$\mu$ & La media del grupo, $\mu=\mathrm{E}[X]$ & $(1+2+6)/3=3$ \\
\midrule
$s$ & Cualquier candidato a representante & 2, 3, 4, 10\dots lo que se quiera probar \\
\midrule
$d_\phi(a,b)$ & El error de reemplazar $a$ por $b$ (Sección 2.3.1) & Con $\phi=x^2$ es $(a-b)^2$ \\
\midrule
$\mathrm{E}[\cdots]$ & El promedio sobre todos los datos del grupo & Se suma el error de cada dato y se divide entre 3 \\
\bottomrule
\end{tabular}
\end{center}
$\mathrm{E}[X]$ es el promedio de los datos y $\mathrm{E}[d(X,s)]$ es el promedio de
los errores: primero se calcula el error de cada dato y luego se promedian.

\subsection{Cada una de las tres partes}

\textbf{Lado izquierdo, $\mathrm{E}[d_\phi(X,s)]$: el error promedio si se elige
$s$.} Con $\phi=x^2$ y $s=2$:
\begin{itemize}
  \item dato 1: $(1-2)^2=1$; dato 2: $(2-2)^2=0$; dato 6: $(6-2)^2=16$;
  \item promedio: $(1+0+16)/3=\mathbf{17/3\approx 5.67}$.
\end{itemize}

\textbf{Primer término de la derecha, $\mathrm{E}[d_\phi(X,\mu)]$: el error
promedio si se elige la media.} Con $\mu=3$:
\begin{itemize}
  \item $(1-3)^2=4$; $(2-3)^2=1$; $(6-3)^2=9$;
  \item promedio: $(4+1+9)/3=\mathbf{14/3\approx 4.67}$, menor que 5.67.
\end{itemize}

\textbf{Segundo término de la derecha, $d_\phi(\mu,s)$: el error de reemplazar la
media por el candidato.} Aquí ya no hay promedio ni datos individuales: es un solo
número, qué tan lejos está el candidato de la media. $(\mu-s)^2=(3-2)^2=\mathbf{1}$.

\subsection{Lo que dice la ecuación}

\[
  \underbrace{\tfrac{17}{3}}_{\text{error con } s=2}
  =\underbrace{\tfrac{14}{3}}_{\text{error con la media}}
  +\underbrace{1}_{\text{lejanía de } s \text{ a la media}}.
\]
\begin{idea}
El error de usar cualquier representante $s$ es el error de usar la media más un
castigo por haberse alejado de ella.
\end{idea}

\subsection{Por qué esto implica que la media es la mejor}

El castigo $d_\phi(\mu,s)$ es una divergencia: nunca es negativa y vale 0 solo
cuando sus dos argumentos son iguales. Si $s=\mu$, el castigo es 0 y el error es
el mínimo posible. Si $s\neq\mu$, el castigo es positivo y el error es mayor.

\begin{center}
\begin{tabular}{ccc}
\toprule
Candidato $s$ & Error promedio & Castigo $(\mu-s)^2$ \\
\midrule
1 & $(0+1+25)/3=8.67$ & 4 \\
2 & $(1+0+16)/3=5.67$ & 1 \\
\textbf{3 (la media)} & $\mathbf{(4+1+9)/3=4.67}$ & \textbf{0} \\
4 & $(9+4+4)/3=5.67$ & 1 \\
5 & $(16+9+1)/3=8.67$ & 4 \\
\bottomrule
\end{tabular}
\end{center}
En cada fila, error $=4.67+$ castigo. La media es la única que no paga castigo.

\subsection{Lo sorprendente: vale para cualquier $\phi$}

Con la I-divergencia los números cambian: el error con la media es $0.750$ y con
$s=2$ es $0.966$. La diferencia es $0.216$, que es $d_I(3,2)=3\ln1.5-1=0.216$. La
forma es idéntica: error con $s$ = error con la media + castigo.

\section{La demostración paso a paso}

(Esta parte es para entender por qué se cumple; no va en la tesis.) Se abre cada
divergencia con la definición
$d_\phi(a,b)=\phi(a)-\phi(b)-\langle a-b,\nabla\phi(b)\rangle$. Solo hace falta
que el promedio es lineal y que las constantes salen del promedio.

\textbf{Paso 1, lado izquierdo.}
$\mathrm{E}[d_\phi(X,s)]=\mathrm{E}\bigl[\phi(X)-\phi(s)-\langle X-s,\nabla\phi(s)\rangle\bigr]$.
Se promedia cada pedazo: $\mathrm{E}[\phi(X)]$ queda igual; $\phi(s)$ es fijo; y
$\mathrm{E}[\langle X-s,\nabla\phi(s)\rangle]=\langle\mu-s,\nabla\phi(s)\rangle$.
\[
  \textbf{(A)}\quad \mathrm{E}[d_\phi(X,s)]=\mathrm{E}[\phi(X)]-\phi(s)-\langle\mu-s,\nabla\phi(s)\rangle .
\]

\textbf{Paso 2, primer término de la derecha.}
$\mathrm{E}[d_\phi(X,\mu)]=\mathrm{E}\bigl[\phi(X)-\phi(\mu)-\langle X-\mu,\nabla\phi(\mu)\rangle\bigr]$.
El último pedazo es $\langle \mathrm{E}[X]-\mu,\nabla\phi(\mu)\rangle=\langle0,\nabla\phi(\mu)\rangle=0$.
Este es el momento clave: desaparece porque el promedio de $X$ es justamente
$\mu$, así que $\mathrm{E}[X-\mu]=0$. Es la única vez que se usa que $\mu$ es la
media.
\[
  \textbf{(B)}\quad \mathrm{E}[d_\phi(X,\mu)]=\mathrm{E}[\phi(X)]-\phi(\mu).
\]

\textbf{Paso 3, segundo término de la derecha.} Sin promedio, solo la definición:
\[
  \textbf{(C)}\quad d_\phi(\mu,s)=\phi(\mu)-\phi(s)-\langle\mu-s,\nabla\phi(s)\rangle .
\]

\textbf{Paso 4.} Al sumar (B) y (C), los términos $-\phi(\mu)$ y $+\phi(\mu)$ se
cancelan y queda $\mathrm{E}[\phi(X)]-\phi(s)-\langle\mu-s,\nabla\phi(s)\rangle$,
que es exactamente (A).

\textbf{Comprobación numérica} ($\phi=x^2$, datos 1, 2, 6, $\mu=3$, $s=2$):
$\mathrm{E}[\phi(X)]=41/3\approx13.67$, $\phi(s)=4$ y
$\langle\mu-s,\nabla\phi(s)\rangle=1\cdot4=4$. Entonces (A) $=13.67-4-4=5.67$,
(B) $=13.67-9=4.67$, (C) $=9-4-4=1$, y (B)+(C) $=5.67=$ (A).

\begin{idea}
El truco es que en (B) el término lineal se hace cero, y eso solo ocurre cuando el
centro del promedio es la media: es el único punto alrededor del cual los desvíos
$X-\mu$ se compensan.
\end{idea}

\section{Consecuencia: K-means generalizado}

El algoritmo alterna dos pasos: \textbf{asignar} cada dato al representante con
menor $d_\phi(x,\text{representante})$ y \textbf{actualizar} cada representante
como la media de su grupo. El segundo paso baja el error (la media es el mejor
representante) y el primero también, por lo que el algoritmo converge para
cualquier divergencia de Bregman. El recíproco (Banerjee et al.) dice que esto
solo ocurre con las divergencias de Bregman: bajo condiciones suaves, ninguna otra
pérdida tiene a la media como predictor óptimo.

\section{Información de Bregman}

\[
  I_\phi(X)=\mathrm{E}\bigl[d_\phi(X,\mathrm{E}[X])\bigr].
\]

\subsection{Qué es}

Es el error que queda \textbf{incluso usando el mejor representante}, la media. Es
el término ``error promedio con la media'' de la ecuación anterior, al que solo se
le da nombre. Mide cuánta variabilidad tiene el grupo. Si todos los datos son
iguales, la media los representa sin error y $I=0$; si están muy dispersos, $I$ es
grande.

\subsection{Con la distancia euclídea es la varianza}

Datos $1$, $2$, $6$, media 3, $\phi=x^2$:
\[
  I_\phi(X)=\frac{(1-3)^2+(2-3)^2+(6-3)^2}{3}=\frac{4+1+9}{3}=\frac{14}{3}\approx4.67.
\]
Ese cálculo es la definición de varianza: el promedio de los desvíos al cuadrado
respecto de la media. Con la I-divergencia, los mismos datos dan $I\approx0.750$:
una ``varianza'' medida con otra regla, que pondera menos las diferencias en
valores grandes.

\subsection{Con KL es la información mutua}

Aquí $X$ no es un número sino una distribución de probabilidad. Hay dos orígenes,
$u_1$ y $u_2$, igual de frecuentes, y un destino $Y$ que puede ser A o B:
\begin{center}
\begin{tabular}{ccc}
\toprule
Origen & $P(A)$ & $P(B)$ \\
\midrule
$u_1$ & 0.9 & 0.1 \\
$u_2$ & 0.1 & 0.9 \\
\bottomrule
\end{tabular}
\end{center}
Cada fila es una distribución condicional $p(Y\mid U=u)$, y esos son los datos $X$,
cada uno con probabilidad $1/2$. La \textbf{media} de esos datos es la
distribución marginal $p(Y)$, el promedio de las filas: $(0.5,0.5)$. El error de
cada fila contra la media, con KL, es
\[
  0.9\ln\frac{0.9}{0.5}+0.1\ln\frac{0.1}{0.5}=0.529-0.161=0.368,
\]
y por simetría lo mismo para $u_2$. El promedio es $I=0.368$, que es la
\textbf{información mutua} $I(U;Y)$: cuánto dice saber el origen sobre el destino.
Se verifica con la fórmula clásica $H(Y)-H(Y\mid U)=0.693-0.325=0.368$. Si los dos
orígenes tuvieran el mismo reparto de destinos, las filas serían iguales a su media
y $I=0$.

\subsection{Por qué importa para la tesis}

\begin{center}
\renewcommand{\arraystretch}{1.3}
\begin{tabular}{@{}>{\raggedright\arraybackslash}p{7.2cm}cc@{}}
\toprule
\textbf{Problema} & $\phi$ & \textbf{Qué se minimiza} \\
\midrule
K-means (agrupamiento por mínima varianza) & $x^2$ & Varianza dentro de grupos \\
\midrule
Co-agrupamiento de Dhillon et al. & KL & Pérdida de información mutua \\
\bottomrule
\end{tabular}
\end{center}
Esa doble lectura justifica usar un solo marco, el de Bregman, para ambos
problemas.

\begin{idea}
La información de Bregman es el error que queda al usar el mejor representante
(la media). Con la distancia euclídea es la varianza; con KL es la información
mutua.
\end{idea}

% =====================================================================
\part{Sección 2.3.3: Relación con las familias exponenciales}\label{sec:parte-expfam}

\section{Por qué existe esta sección}

Hasta ahora se dijo que la I-divergencia es la apropiada para conteos, pero eso es
un argumento de sentido común: pondera el error según el tamaño de la celda. Un
jurado puede preguntar por qué esa y no otra función que también pese por tamaño.
La respuesta estadística de esta sección es:

\begin{idea}
Minimizar la I-divergencia es exactamente lo mismo que maximizar la verosimilitud
de un modelo de Poisson.
\end{idea}

Eso cambia el estatus del método: deja de ser una regla de ajuste elegida a ojo y
pasa a ser el ajuste de un modelo estadístico con un supuesto claro, que los
conteos de cada celda son Poisson. Tiene tres consecuencias:
\begin{itemize}
  \item \textbf{Justificación:} se elige la I-divergencia porque corresponde al
  modelo natural para conteos.
  \item \textbf{Lectura de los resultados:} un bloque deja de ser ``un grupo de
  celdas parecidas'' y es una región donde los viajes se explican con una
  intensidad media común. Es lo que se usa en la Sección 2.5 y en la descomposición
  de la divergencia.
  \item \textbf{Conexión con el capítulo de estadística:} el modelo de Poisson con
  interacción (log-lineal) es el marco para decir qué es ``interacción
  espacio-temporal''.
\end{itemize}

\section{Qué es una familia exponencial}

Es una colección de distribuciones que comparten una forma matemática: la normal,
la Poisson, la binomial, la gamma. No hace falta la forma general. Lo que importa
es que \textbf{cada una tiene un parámetro de media $\mu$}, el valor promedio de
los datos, y que la probabilidad de un dato depende de qué tan cerca está de esa
media.

\section{El resultado de Banerjee et al.\ (2005)}

Existe una \textbf{correspondencia uno a uno} entre divergencias de Bregman y
familias exponenciales. Para cada familia hay una $\phi$, y la densidad se escribe
\[
  p(x\mid\mu)=\exp\{-d_\phi(x,\mu)\}\;b_\phi(x).
\]
\begin{center}
\renewcommand{\arraystretch}{1.3}
\begin{tabular}{@{}>{\raggedright\arraybackslash}p{3.2cm}>{\raggedright\arraybackslash}p{10.6cm}@{}}
\toprule
\textbf{Símbolo} & \textbf{Qué es} \\
\midrule
$p(x\mid\mu)$ & Probabilidad de observar $x$ si la media verdadera es $\mu$ \\
\midrule
$d_\phi(x,\mu)$ & La divergencia entre el dato y la media \\
\midrule
$\exp\{-d\}$ & Una función que baja cuando $d$ sube: más lejos de la media, menos probable \\
\midrule
$b_\phi(x)$ & Un factor que depende solo del dato, no de $\mu$ \\
\bottomrule
\end{tabular}
\end{center}
La lectura es directa: un dato es probable cuando está cerca de la media, y
``cerca'' se mide con la divergencia de Bregman de esa familia. Cada familia
tiene su propia regla de cercanía:
\begin{center}
\begin{tabular}{lcc}
\toprule
Familia & $\phi$ & Divergencia \\
\midrule
Normal (varianza $\sigma^2$ fija) & $x^2/(2\sigma^2)$ & $(x-\mu)^2/(2\sigma^2)$, la euclídea escalada \\
\midrule
Poisson & $x\log x-x$ & la I-divergencia \\
\bottomrule
\end{tabular}
\end{center}
Esto prolonga la idea de la Sección 2.3.1, ``cada $\phi$ corresponde a un supuesto
sobre cómo varía el error'': ahora el supuesto tiene nombre. La euclídea es el
error normal y la I-divergencia es el error Poisson.

\section{Verificación para Poisson}

Si $x\sim\mathrm{Poisson}(\mu)$, la probabilidad es $p(x\mid\mu)=e^{-\mu}\mu^x/x!$
y, tomando logaritmo,
\[
  \log p(x\mid\mu)=x\log\mu-\mu-\log x!.
\]
La I-divergencia desarrollada es $d_I(x,\mu)=x\log x-x\log\mu-x+\mu$, así que
$-d_I(x,\mu)=-x\log x+x\log\mu+x-\mu$. Sumando y restando $x\log x-x-\log x!$ se
obtiene la fórmula de Poisson:
\[
  \log p(x\mid\mu)=-d_I(x,\mu)+\underbrace{\bigl(x\log x-x-\log x!\bigr)}_{\text{no depende de }\mu}.
\]
El último paréntesis solo depende del dato, así que no influye al elegir $\mu$.

\textbf{Con números}, $x=4$ y $\mu=3$:
\begin{itemize}
  \item Probabilidad: $e^{-3}\cdot3^4/4!=0.168$, y $\log p=-1.784$.
  \item $d_I(4,3)=4\ln(4/3)-1=0.151$.
  \item Parte constante: $4\ln4-4-\ln24=-1.633$.
  \item $-0.151-1.633=-1.784$. \checkmark
\end{itemize}

\section{Por qué minimizar la divergencia equivale a maximizar la verosimilitud}

En la fórmula anterior, para elegir el mejor $\mu$ el término constante no
importa; lo que cambia con $\mu$ es solo $-d_I(x,\mu)$. Mayor verosimilitud
significa mayor $\log p$, es decir, mayor $-d_I$, es decir, menor $d_I$. Con
$x=4$ y dos medias candidatas:
\begin{center}
\begin{tabular}{ccc}
\toprule
$\mu$ & $d_I(4,\mu)$ & $\log p$ \\
\midrule
3 & 0.151 & $-1.784$ \\
5 & 0.107 & $-1.740$ \\
\bottomrule
\end{tabular}
\end{center}
$\mu=5$ tiene menor divergencia y mayor verosimilitud: el orden es el mismo en las
dos columnas, siempre.

\section{Para una matriz completa}

Si cada celda es un conteo Poisson independiente, $O_{ij}\sim\mathrm{Poisson}(\mu_{ij})$,
el logaritmo de la probabilidad conjunta es la suma de las celdas. Como la
I-divergencia es aditiva por celda,
\[
  D_I(O\,\|\,\mu)=\sum_{i,j}d_I(O_{ij},\mu_{ij})=\tfrac12\,\mathrm{Dev}(\mu),
\]
donde $\mathrm{Dev}(\mu)$ es la \textbf{devianza} del modelo de Poisson, la medida
estándar de qué tan mal ajusta un modelo a los conteos (Agresti, 2013). Aquí
$\mu_{ij}$ es lo que predice el modelo para cada celda, en esta tesis la media del
bloque al que pertenece; $D_I$ mide el desajuste entre los conteos reales y esa
predicción. Elegir los bloques que minimizan $D_I$ es elegir el modelo de Poisson
por bloques que mejor explica los datos.

\begin{idea}
La I-divergencia no es una elección arbitraria: es la divergencia que corresponde
al modelo de Poisson. Por eso el co-agrupamiento con I-divergencia se puede leer
como el ajuste de un modelo estadístico a los conteos.
\end{idea}

% =====================================================================
\part{Sección 2.4: Co-agrupamiento de Bregman}

\section{Por qué existe esta sección}

Hasta ahora hay dos piezas separadas: las \textbf{divergencias de Bregman} (2.3),
una forma de medir el error al reemplazar un dato por un representante, y el
\textbf{co-agrupamiento} (2.2), que agrupa a la vez las filas (regiones) y las
columnas (intervalos de tiempo) y forma un tablero de bloques. La pregunta que las
une es: ¿cómo se reconstruye la matriz a partir de los bloques y cómo se mide qué
tan buena es esa reconstrucción?

\begin{idea}
Se resume la matriz con unos pocos promedios, se reconstruye una matriz aproximada
a partir de ellos y se eligen los grupos que hacen mínimo el error de la
aproximación, medido con una divergencia de Bregman.
\end{idea}

Es la sección que convierte la teoría de la 2.3 en un método concreto, el que luego
usan ST-CC y la tesis.

\begin{center}
\renewcommand{\arraystretch}{1.3}
\begin{tabular}{@{}>{\raggedright\arraybackslash}p{4.2cm}>{\raggedright\arraybackslash}p{9.6cm}@{}}
\toprule
\textbf{Subsección} & \textbf{Pregunta que responde} \\
\midrule
2.4.1 Notación & ¿Cómo se escribe un co-agrupamiento? \\
\midrule
2.4.2 Aproximación y MBI & ¿Qué promedios se conservan y cómo se reconstruye la matriz? \\
\midrule
2.4.3 Las seis bases & ¿Qué opciones hay para elegir los promedios? \\
\midrule
2.4.4 Algoritmo & ¿Cómo se encuentran los grupos? \\
\bottomrule
\end{tabular}
\end{center}
Estas notas cubren 2.4.1 y 2.4.2, que son las que más cuestan.

\section{Ejemplo que se usa}

Una matriz de 4 regiones por 4 intervalos, agrupadas de dos en dos:
\[
  O=\begin{pmatrix}
    12 & 8 & 3 & 1\\
    8 & 12 & 1 & 3\\
    2 & 2 & 10 & 6\\
    2 & 2 & 6 & 10
  \end{pmatrix}.
\]
Las regiones $r_1,r_2$ forman el grupo 1 de filas y $r_3,r_4$ el grupo 2. Los
intervalos $t_1,t_2$ forman el grupo 1 de columnas y $t_3,t_4$ el grupo 2. Es un
tablero de $2\times2$ bloques.

\section{2.4.1 Notación}

\subsection{Las funciones $\rho$ y $\tau$}

$\rho$ asigna cada fila a un grupo de filas: $\rho(r_1)=\rho(r_2)=1$ y
$\rho(r_3)=\rho(r_4)=2$. $\tau$ hace lo mismo con las columnas. Un co-agrupamiento
es el par $(\rho,\tau)$, y $K_R$ y $K_T$ son cuántos grupos hay de cada tipo (aquí
2 y 2).

\subsection{Los bloques}

$I_g$ es el conjunto de filas del grupo $g$, $J_h$ el de columnas del grupo $h$, y
$B_{gh}=I_g\times J_h$ es el bloque que forman. Por ejemplo, $B_{11}$ son las
cuatro celdas de arriba a la izquierda.

\subsection{La lectura probabilística}

Para escribir los promedios de forma compacta, Banerjee et al.\ eligen una celda
de la matriz \textbf{al azar} y llaman $U$ a la fila elegida y $V$ a la columna
elegida (cada una uniforme), $Z=O_{UV}$ al valor de esa celda (una variable
aleatoria), y $\hat U=\rho(U)$, $\hat V=\tau(V)$ a los grupos de la fila y la
columna elegidas. Con eso, cada promedio es una \emph{esperanza condicional}: el
promedio de $Z$ restringido a lo que se indica.

\begin{center}
\renewcommand{\arraystretch}{1.3}
\begin{tabular}{@{}>{\raggedright\arraybackslash}p{4.4cm}>{\raggedright\arraybackslash}p{5.6cm}>{\raggedright\arraybackslash}p{3.4cm}@{}}
\toprule
\textbf{Notación} & \textbf{Significado} & \textbf{En el ejemplo} \\
\midrule
$\mathrm{E}[Z]$ & Promedio de toda la matriz & $88/16=5.5$ \\
\midrule
$\mathrm{E}[Z\mid U=i]$ & Media de la fila $i$ & Fila $r_1$: $24/4=6$ \\
\midrule
$\mathrm{E}[Z\mid\hat U=g,\hat V=h]$ & Media del bloque $B_{gh}$ & $B_{11}$: $40/4=10$ \\
\bottomrule
\end{tabular}
\end{center}
No hay nada profundo aquí: es una forma de escribir ``promedio de las celdas que
cumplen esta condición'' con una sola notación. La ventaja es que la información de
Bregman y la media como mejor representante (2.3.2) se aplican directamente.

Medias de los cuatro bloques del ejemplo:
\begin{center}
\begin{tabular}{ccc}
\toprule
Bloque & Celdas & Media \\
\midrule
$B_{11}$ & 12, 8, 8, 12 & 10 \\
$B_{12}$ & 3, 1, 1, 3 & 2 \\
$B_{21}$ & 2, 2, 2, 2 & 2 \\
$B_{22}$ & 10, 6, 6, 10 & 8 \\
\bottomrule
\end{tabular}
\end{center}

\section{2.4.2 Aproximación y principio de mínima información de Bregman}

\subsection{La pregunta}

Un co-agrupamiento resume la matriz con un conjunto de promedios. Hay que decidir
cuáles promedios conservar y cómo reconstruir la matriz a partir de ellos. Banerjee
et al.\ responden con dos elementos.

\subsection{Elemento 1: la base $\mathcal C$}

Es la lista de promedios que se conservan. Cada elemento es un conjunto de
variables tomadas de $\{U,V,\hat U,\hat V\}$. La base $\{\hat U,\hat V\}$ conserva
la media de cada bloque (es la más simple, C2). La base $\{\hat U\},\{\hat V\}$
conserva solo la media de cada grupo de filas y la de cada grupo de columnas, sin
saber qué pasa en cada bloque (C1). Hay seis bases en total, que se ven en 2.4.3.

\subsection{Elemento 2: el principio MBI}

Con los promedios fijados hay muchas matrices que los respetan, y hay que elegir
una. La regla es: entre todas las matrices que conservan los promedios de la base,
se elige la de \textbf{menor información de Bregman}. Recuerda de la 2.3.2 que la
información de Bregman mide cuánta variabilidad hay. Elegir la de menor
información es elegir la matriz más simple compatible con los promedios: no
inventa estructura que los promedios no exigen. Es el mismo espíritu que la máxima
entropía, y con la I-divergencia o KL coinciden.

\subsection{Cómo se ve en el ejemplo}

\textbf{Base C2} (media de cada bloque): la matriz más simple compatible rellena
cada bloque con su media,
\[
  \hat O=\begin{pmatrix}
    10 & 10 & 2 & 2\\
    10 & 10 & 2 & 2\\
    2 & 2 & 8 & 8\\
    2 & 2 & 8 & 8
  \end{pmatrix}.
\]

\textbf{Base C1} (solo media por grupo de filas y por grupo de columnas): el MBI
con I-divergencia da la fórmula multiplicativa
$\mathrm{E}[Z\mid\hat U]\,\mathrm{E}[Z\mid\hat V]/\mathrm{E}[Z]$. Las medias por
grupo de filas son 6 y 5, las de columnas son 6 y 5, y $\mathrm{E}[Z]=5.5$:
\begin{center}
\begin{tabular}{ccc}
\toprule
Bloque & Cálculo & Aproximación \\
\midrule
$B_{11}$ & $6\cdot6/5.5$ & 6.55 \\
$B_{12}$ & $6\cdot5/5.5$ & 5.45 \\
$B_{21}$ & $5\cdot6/5.5$ & 5.45 \\
$B_{22}$ & $5\cdot5/5.5$ & 4.55 \\
\bottomrule
\end{tabular}
\end{center}
Es el modelo de independencia: las medias de fila y de columna se multiplican. Esta
matriz respeta las medias por grupo (la media del grupo 1 de filas es
$(6.55+5.45)/2=6$) pero no ve que el bloque $B_{11}$ es mucho más alto que
$B_{12}$. La base elige qué estructura puede capturar la aproximación: C2 captura
diferencias entre bloques; C1 solo captura efectos de fila y de columna por
separado.

\subsection{Propiedades que se mencionan, sin demostrar}

La solución MBI, $\hat Z_{\mathcal C}$, es además la mejor aproximación en
divergencia esperada entre las de su misma forma, y cumple una propiedad
pitagórica que se usa en la Sección 2.5 para descomponer la divergencia.

\section{La función objetivo}

\[
  (\rho^*,\tau^*)=\arg\min_{\rho,\tau}\;\mathrm{E}\bigl[d_\phi\bigl(Z,\hat Z_{\mathcal C}(\rho,\tau)\bigr)\bigr].
\]
Se lee: entre todos los co-agrupamientos posibles, se elige el que hace mínimo el
error promedio entre la matriz real y su aproximación. Con medida uniforme es
proporcional a $\sum_{i,j}d_\phi(O_{ij},\hat O_{ij})$, la suma de errores celda por
celda, como en la 2.3.1.

\textbf{Con números} (base C2, I-divergencia, los grupos del ejemplo):
\begin{center}
\begin{tabular}{ccc}
\toprule
Bloque & Media & Suma de $d_I$ celda por celda \\
\midrule
$B_{11}$ & 10 & 0.805 \\
$B_{12}$ & 2 & 1.047 \\
$B_{21}$ & 2 & 0 (todas valen 2) \\
$B_{22}$ & 8 & 1.011 \\
\midrule
\textbf{Total} & & $\approx 2.86$ \\
\bottomrule
\end{tabular}
\end{center}
El algoritmo busca la partición que haga ese total lo más chico posible.

\begin{idea}
Un co-agrupamiento resume la matriz con promedios (la base). El MBI reconstruye la
matriz más simple que respeta esos promedios. Se elige la partición cuya
reconstrucción se parece más a los datos, medido con una divergencia de Bregman.
\end{idea}

\section{2.4.3 Las seis bases}

\subsection{Por qué existe esta subsección}

En la 2.4.2 se vio que hay que decidir qué promedios conservar (la base), y que
cada opción produce una aproximación distinta. Banerjee et al.\ demuestran que en
dos dimensiones solo hay \textbf{seis bases distintas}: las demás son triviales (no
resumen nada) o redundantes (equivalen a alguna de estas). Cada base responde a la
pregunta: además de dividir en bloques, ¿cuánto detalle de la matriz se quiere
conservar?

\subsection{Cómo leer la notación}

$U$ y $V$ son una fila y una columna \textbf{individuales}; $\hat U$ y $\hat V$
son el \textbf{grupo} de esa fila y de esa columna. Un conjunto como
$\{\hat U,\hat V\}$ significa ``conserva la media de cada pareja (grupo de filas,
grupo de columnas)'', es decir, cada bloque; $\{U\}$ significa ``conserva la media
de cada fila individual''.

\subsection{Las seis bases en palabras}

\begin{center}
\renewcommand{\arraystretch}{1.3}
\footnotesize
\begin{tabular}{@{}l>{\raggedright\arraybackslash}p{4.4cm}>{\raggedright\arraybackslash}p{5.6cm}>{\raggedright\arraybackslash}p{2.4cm}@{}}
\toprule
\textbf{Base} & \textbf{Qué conserva} & \textbf{Qué dice la aproximación} & \textbf{Familia} \\
\midrule
C1 & Media por grupo de filas y por grupo de columnas & Independencia entre filas y columnas, sin bloques & Sin bloques \\
\midrule
C2 & Media de cada bloque & Todas las celdas de un bloque valen lo mismo & Bloques \\
\midrule
C3 & Bloque y cada fila & Bloque, corregido por cuánto más grande es la fila que su grupo & Bloques y volumen de fila \\
\midrule
C4 & Bloque y cada columna & Bloque, corregido por el volumen de la columna & Bloques y volumen de columna \\
\midrule
C5 & Bloque, cada fila y cada columna & Bloque, corregido por el volumen de la fila y de la columna & Bloques y ambos volúmenes \\
\midrule
C6 & Cada fila dentro de cada grupo de columnas, y cada columna dentro de cada grupo de filas & Perfiles internos de cada bloque & Más fino \\
\bottomrule
\end{tabular}
\end{center}

\subsection{La idea de las fórmulas con I-divergencia}

La más fácil de leer es C5:
\[
  \hat O_{ij}=\frac{\mathrm{E}[Z|U]\,\mathrm{E}[Z|V]\,\mathrm{E}[Z|\hat U,\hat V]}
  {\mathrm{E}[Z|\hat U]\,\mathrm{E}[Z|\hat V]}.
\]
Reagrupada:
\[
  \hat O_{ij}=\underbrace{\mathrm{E}[Z|\hat U,\hat V]}_{\text{media del bloque}}
  \times\underbrace{\frac{\mathrm{E}[Z|U]}{\mathrm{E}[Z|\hat U]}}_{\text{factor de fila}}
  \times\underbrace{\frac{\mathrm{E}[Z|V]}{\mathrm{E}[Z|\hat V]}}_{\text{factor de columna}}.
\]
Cada factor compara una fila (o columna) con \emph{su grupo}: si una región tiene
el doble de viajes que el promedio de su grupo, el factor de fila es 2. La
aproximación toma la media del bloque y la escala por esos dos factores. Por eso
C5 distingue entre ``esta región es más grande'' y ``esta región tiene un patrón
distinto''. C3 aplica solo el factor de fila, C4 solo el de columna y C2 ninguno.

\subsection{Ejemplo: la diferencia entre C2 y C5}

Para verla hace falta una matriz donde las filas tengan distinto volumen. Una
matriz con un solo bloque (todo en un grupo):
\[
  \begin{pmatrix}2&4\\6&12\end{pmatrix}.
\]
La fila 2 tiene el triple de viajes que la fila 1, con el mismo perfil. Las medias
de fila son 3 y 9, las de columna 4 y 8, y la global 6.

\textbf{Con C2} (media del bloque) todas las celdas valen 6. La I-divergencia
suma $d(2,6)=1.80$, $d(4,6)=0.38$, $d(6,6)=0$ y $d(12,6)=2.32$, es decir,
$\approx\mathbf{4.50}$.

\textbf{Con C5}, con un solo bloque la fórmula queda
$\mathrm{E}[Z|U]\,\mathrm{E}[Z|V]/\mathrm{E}[Z]$. Celda $(1,1)$: $3\cdot4/6=2$;
$(1,2)$: $3\cdot8/6=4$; $(2,1)$: $9\cdot4/6=6$; $(2,2)$: $9\cdot8/6=12$.
Reproduce la matriz exacta: pérdida \textbf{0}.

La matriz es un patrón único multiplicado por el volumen de cada fila y de cada
columna. C2 lo ve como un error grande; C5 lo explica entero porque descuenta el
volumen. Esta es la diferencia que la tesis usa para justificar C5: las
diferencias de volumen entre regiones no son interacción espacio-temporal.

\subsection{Casos particulares conocidos}

\begin{center}
\renewcommand{\arraystretch}{1.3}
\begin{tabular}{@{}>{\raggedright\arraybackslash}p{4.4cm}>{\raggedright\arraybackslash}p{9.4cm}@{}}
\toprule
\textbf{Base y divergencia} & \textbf{Método conocido} \\
\midrule
C2 + euclídea & Co-agrupamiento clásico de bloques constantes (primer residuo de Cho et al., 2004) \\
\midrule
C6 + euclídea & Residuo cuadrático mínimo de Cho et al. (segundo residuo); sobre él se construye el BCC\_MSSR de Wu et al.\ (2022) \\
\midrule
C5 + I-divergencia & ITCC de Dhillon et al.\ (2003): con distribuciones conjuntas la I-divergencia es KL, y minimiza la pérdida de información mutua \\
\bottomrule
\end{tabular}
\end{center}
Esto respalda la idea de la 2.3: un solo marco generaliza métodos que parecían
distintos.

\subsection{La nota sobre C6}

C6 aparece en dos contextos con la misma estructura y distinta combinación. Con la
euclídea combina las marginales del bloque \emph{sumando} (el residuo de Cho). Con
la I-divergencia las combina \emph{multiplicando} (la prueba de independencia
dentro del bloque de la Sección 2.5). Es la misma idea con dos álgebras: la
euclídea trabaja con diferencias y la I-divergencia con razones, igual que en la
2.3.1 la euclídea mide diferencias y la I-divergencia pondera por tamaño.

\begin{idea}
Las seis bases son seis niveles de detalle para reconstruir la matriz. C2 conserva
solo bloques; C5 añade el volumen de cada fila y columna; C6 conserva más
estructura interna. Cada una, con una divergencia, reproduce un método clásico.
\end{idea}

\section{2.4.4 Algoritmo}

\subsection{Por qué existe esta subsección}

Ya se sabe qué se quiere: la partición $(\rho,\tau)$ que haga mínimo el error entre
la matriz y su aproximación. Falta \textbf{cómo encontrarla}. La tentación es
probar todas las particiones, pero es inviable: solo asignar 100 filas a 5 grupos
admite unas $5^{100}\approx10^{70}$ posibilidades, sin contar las columnas.

El texto lo dice con un término técnico: el problema es \textbf{NP-difícil}, por
reducción desde K-means. Significa que no se conoce un método eficiente que
garantice la mejor partición. Por eso Banerjee et al.\ proponen un método que no
garantiza el óptimo pero llega a una buena solución rápido: la
\textbf{minimización alternante}.

\subsection{La idea}

Es la misma que K-means (2.3.2): alternar entre dos pasos sencillos. Con los
\textbf{grupos fijos}, calcular los mejores promedios (los de la base); con los
\textbf{promedios fijos}, mejorar los grupos. Aquí hay dos tipos de grupos (filas
y columnas), así que se alterna entre tres cosas: promedios, filas y columnas.

\subsection{Los pasos}

\begin{center}
\renewcommand{\arraystretch}{1.3}
\begin{tabular}{@{}c>{\raggedright\arraybackslash}p{6.4cm}>{\raggedright\arraybackslash}p{6.4cm}@{}}
\toprule
\textbf{Paso} & \textbf{Qué hace} & \textbf{En palabras} \\
\midrule
1 & Inicializa $\rho$ y $\tau$ al azar & Empieza con grupos cualesquiera \\
\midrule
2 & Calcula los estadísticos de la base y la aproximación $\hat O$ & Con los grupos actuales, calcula los promedios y reconstruye la matriz \\
\midrule
3 & Reasigna cada fila al grupo $g$ que minimiza $\sum_j d_\phi(O_{ij},\hat O_{ij})$ calculada como si $i$ perteneciera a $g$; recalcula los estadísticos & Cada fila se muda al grupo que mejor la explica \\
\midrule
4 & Repite el paso 3 con las columnas, con $\rho$ fija & Lo mismo para las columnas \\
\midrule
5 & Si la reducción de la divergencia es menor que $\varepsilon$, termina; si no, vuelve al paso 3 & Para cuando casi no mejora \\
\bottomrule
\end{tabular}
\end{center}

El paso 3 es el que cuesta. Para decidir dónde va la fila $i$, se prueba cada
grupo $g$ y se calcula el error de la fila si estuviera ahí, usando los promedios
de ese grupo. Se elige el grupo de menor error. Es el equivalente a ``asignar cada
punto al centroide más cercano'' de K-means, solo que el ``centroide'' ahora es un
conjunto de promedios por bloque.

\subsection{Ejemplo con la matriz 4$\times$4}

Matriz de la 2.4.2, base C2 e I-divergencia:
\[
  O=\begin{pmatrix}
    12 & 8 & 3 & 1\\
    8 & 12 & 1 & 3\\
    2 & 2 & 10 & 6\\
    2 & 2 & 6 & 10
  \end{pmatrix}.
\]

\textbf{Paso 1.} Se supone que $\tau$ ya es correcta (columnas 1, 2 y 3, 4) y que
$\rho$ empieza mal: las filas $r_1,r_2,r_3$ en el grupo 1 y $r_4$ sola en el
grupo 2. El error total de esa partición es $\approx\mathbf{15.6}$.

\textbf{Paso 2.} Promedios de los bloques:
\begin{center}
\begin{tabular}{ccc}
\toprule
Bloque & Celdas & Media \\
\midrule
Grupo 1, columnas 1--2 & 12, 8, 8, 12, 2, 2 & 7.33 \\
Grupo 1, columnas 3--4 & 3, 1, 1, 3, 10, 6 & 4 \\
Grupo 2, columnas 1--2 & 2, 2 & 2 \\
Grupo 2, columnas 3--4 & 6, 10 & 8 \\
\bottomrule
\end{tabular}
\end{center}

\textbf{Paso 3.} Se mira la fila $r_3=(2,2,10,6)$. Si se queda en el grupo 1, se
la compara con $(7.33,7.33,4,4)$ y el error es $2.73+2.73+3.16+0.43\approx
\mathbf{9.07}$. Si pasa al grupo 2, se la compara con $(2,2,8,8)$ y el error es
$0+0+0.23+0.27\approx\mathbf{0.51}$. Se muda al grupo 2. Las filas $r_1$ y $r_2$
prefieren claramente el grupo 1 (error $\approx3$ frente a $\approx24$) y $r_4$ se
queda donde está. La partición de filas queda $\{r_1,r_2\},\{r_3,r_4\}$: la
correcta, con error total $\approx\mathbf{2.86}$.

\textbf{Paso 4.} Con $\rho$ ya correcta, ninguna columna mejora cambiando de grupo.

\textbf{Paso 5.} El error bajó de 15.6 a 2.86 en el paso 3 y en el paso 4 no
cambió. La mejora es menor que $\varepsilon$ y el algoritmo termina.

\subsection{Por qué converge}

\begin{itemize}
  \item \textbf{Cada paso no aumenta el error.} Al recalcular los promedios, la
  media es el mejor representante de cada grupo (esto es la 2.3.2). Al reasignar
  filas, cada una elige el grupo que menos error le da.
  \item \textbf{Hay un número finito de particiones.} Como el error nunca sube y
  no hay infinitos estados, el algoritmo no puede dar vueltas para siempre: se
  detiene en un número finito de iteraciones.
\end{itemize}

\subsection{La limitación: mínimo local}

Converger no significa llegar a la mejor solución. El algoritmo termina en un
\textbf{mínimo local}: una partición que no mejora cambiando las filas, las
columnas o los promedios, pero que podría ser peor que otra partición lejana. Es
como un excursionista que baja una montaña por el camino más empinado: llega a un
valle, pero puede haber otro valle más profundo al otro lado.

Tres precauciones que menciona el texto:
\begin{enumerate}
  \item \textbf{Varias inicializaciones.} Se ejecuta con muchos arranques al azar
  y se conserva el mejor resultado. Liu et al.\ usaron 200.
  \item \textbf{Búsqueda local.} Cho et al.\ muestran que, después de converger,
  mover filas o columnas de una en una puede seguir bajando el error. El algoritmo
  básico mueve todas a la vez (por lotes), y por eso hay mejoras que se le
  escapan.
  \item \textbf{Costo.} Cada iteración cuesta $O\bigl(mn(K_R+K_T)\bigr)$
  operaciones: proporcional al tamaño de la matriz por el número de grupos.
\end{enumerate}

\begin{idea}
Como probar todas las particiones es imposible, se alterna entre calcular
promedios y reasignar filas y columnas. Cada paso baja el error, así que converge,
pero a un mínimo local, por lo que se repite con muchos arranques.
\end{idea}

% =====================================================================
\part{Sección 2.5: La I-divergencia y el modelo de Poisson}

\section{Por qué existe esta sección}

El marco de Bregman (2.4) funciona con cualquier divergencia. Pero la tesis usa una
sola, la I-divergencia, con conteos de viajes. Esta sección particulariza todo a
ese caso y obtiene tres cosas que solo existen ahí:
\begin{enumerate}
  \item \textbf{Fórmulas explícitas} para las bases C2, C5 y C6 con totales de
  fila, columna y bloque.
  \item \textbf{Lectura estadística:} cada base es un modelo de Poisson concreto,
  y tres de ellas están anidadas (C2 $\subset$ C5 $\subset$ C6).
  \item \textbf{Una descomposición} de la divergencia en tres términos con
  interpretación propia, que justifica el estadístico de heterogeneidad $G^2$ del
  método propuesto.
\end{enumerate}

\begin{idea}
La 2.4 dice cómo se hace el co-agrupamiento; la 2.5 dice qué significa lo que se
mide y por qué $G^2$ es la herramienta adecuada para buscar patrones locales.
\end{idea}

\begin{center}
\renewcommand{\arraystretch}{1.3}
\begin{tabular}{@{}>{\raggedright\arraybackslash}p{5.4cm}>{\raggedright\arraybackslash}p{8.4cm}@{}}
\toprule
\textbf{Subsección} & \textbf{Pregunta} \\
\midrule
2.5.1 Tres bases en notación matricial & ¿Cómo se calculan C2, C5 y C6 con totales? \\
\midrule
2.5.2 Modelos log-lineales de Poisson & ¿Qué modelo estadístico es cada base? \\
\midrule
2.5.3 Descomposición & ¿De qué partes se compone el error? \\
\midrule
2.5.4 Escala & ¿Importa si se usan conteos o promedios? \\
\bottomrule
\end{tabular}
\end{center}
Estas notas cubren toda la sección 2.5.

\section{2.5.1 Tres bases en notación matricial}

\subsection{Los totales que se necesitan}

Con la matriz de conteos $O$ de 4$\times$4 de la 2.4.2:
\begin{center}
\renewcommand{\arraystretch}{1.3}
\begin{tabular}{@{}>{\raggedright\arraybackslash}p{2.8cm}>{\raggedright\arraybackslash}p{7.0cm}>{\raggedright\arraybackslash}p{3.8cm}@{}}
\toprule
\textbf{Símbolo} & \textbf{Significado} & \textbf{Ejemplo} \\
\midrule
$O_{i\cdot}$ & Total de la fila $i$ & $r_1=24$, $r_3=20$ \\
\midrule
$O_{\cdot j}$ & Total de la columna $j$ & $c_1=24$, $c_3=20$ \\
\midrule
$N$ & Total general & 88 \\
\midrule
$O_{g\cdot},O_{\cdot h}$ & Total del grupo de filas $g$ y del grupo de columnas $h$ & 48 y 40 \\
\midrule
$N_{gh}$ & Total del bloque & 40, 8, 8, 32 \\
\midrule
$O_{ih}$ & Total de la fila $i$ \emph{dentro} del grupo de columnas $h$ & $r_1$ en columnas 1--2: 20 \\
\midrule
$O_{gj}$ & Total de la columna $j$ \emph{dentro} del grupo de filas $g$ & $c_1$ en filas 1--2: 20 \\
\bottomrule
\end{tabular}
\end{center}
Son los mismos promedios de la 2.4, escritos como sumas (promedio por número de
celdas).

\subsection{Las tres fórmulas}

\[
  \hat O^{\mathrm{C2}}_{ij}=\frac{N_{gh}}{|I_g||J_h|},\qquad
  \hat O^{\mathrm{C5}}_{ij}=\frac{O_{i\cdot}O_{\cdot j}N_{gh}}{O_{g\cdot}O_{\cdot h}},\qquad
  \hat O^{\mathrm{C6}}_{ij}=\frac{O_{ih}O_{gj}}{N_{gh}}.
\]
\begin{center}
\renewcommand{\arraystretch}{1.3}
\begin{tabular}{@{}l>{\raggedright\arraybackslash}p{12.4cm}@{}}
\toprule
\textbf{Base} & \textbf{Lectura} \\
\midrule
C2 & Todas las celdas del bloque valen la media del bloque \\
\midrule
C5 & Reparte el total del bloque en proporción al volumen global de cada región y cada intervalo \\
\midrule
C6 & Trata cada bloque como una tabla de contingencia \emph{independiente}: lo esperado si filas y columnas no estuvieran asociadas \emph{dentro del bloque} \\
\bottomrule
\end{tabular}
\end{center}
Con el ejemplo $\begin{pmatrix}2&4\\6&12\end{pmatrix}$ de un solo bloque, C5 da en
la celda $(1,1)$: $6\cdot8\cdot24/(24\cdot24)=2$, el dato exacto.

\subsection{C5 y la información mutua}

\[
  D_I(O\,\|\,\hat O^{\mathrm{C5}})=N\,\bigl[I(R;T)-I(\hat R;\hat T)\bigr].
\]
Con C5, el error de la aproximación es \textbf{la información mutua que se pierde}
al agrupar las regiones y los intervalos ($R$, $T$ son filas y columnas; $\hat R$,
$\hat T$ sus grupos), multiplicada por $N$. Es la conexión con Dhillon (ITCC) vista
en la 2.3: la información de Bregman con KL es información mutua.

\subsection{C6 y la prueba de independencia}

\[
  D_I(O\,\|\,\hat O^{\mathrm{C6}})=\tfrac12\sum_{g,h}G^2_{gh},\qquad
  G^2_{gh}=2\sum_{(i,j)\in B_{gh}}O_{ij}\log\frac{O_{ij}N_{gh}}{O_{ih}O_{gj}}.
\]
$G^2$ es el estadístico de razón de verosimilitud de la \textbf{prueba de
independencia} de una tabla de contingencia, el que se usa en estadística clásica
para ver si dos variables están asociadas. También se escribe
$G^2_{gh}=2N_{gh}\hat I_{gh}(R;T)$: la información mutua dentro del bloque,
escalada por su masa.

\textbf{Ejemplo con un bloque con patrón.} Sea el bloque
$\begin{pmatrix}6&2\\2&6\end{pmatrix}$. Los totales de fila y de columna son 8 y
$N=16$, así que lo esperado bajo independencia es $8\cdot8/16=4$ en cada celda.
Pero los datos son 6, 2, 2, 6: hay asociación (la diagonal pesa más).
\[
  G^2=2\bigl[2\cdot6\ln1.5+2\cdot2\ln0.5\bigr]=2[4.866-2.773]\approx\mathbf{4.19}.
\]
Un bloque sin asociación (todas las celdas valen 4) daría $G^2=0$. Se verifica con
la I-divergencia: $d(6,4)=0.433$ y $d(2,4)=0.614$, y $2(0.433+0.614)=2.093=G^2/2$.

\begin{idea}
$G^2$ mide cuánta asociación entre regiones y horarios queda dentro del bloque.
Este es el patrón local que busca la tesis.
\end{idea}

\section{2.5.2 Equivalencia con modelos log-lineales de Poisson}

\subsection{De dónde viene}

En la 2.3.3 se vio que minimizar la I-divergencia es maximizar la verosimilitud de
Poisson. Falta saber \textbf{qué modelo} de Poisson. Cada base define uno, para las
medias $\mu_{ij}$ de las celdas:
\begin{center}
\renewcommand{\arraystretch}{1.3}
\begin{tabular}{@{}l>{\raggedright\arraybackslash}p{5.2cm}>{\raggedright\arraybackslash}p{6.4cm}@{}}
\toprule
\textbf{Base} & \textbf{Modelo} & \textbf{Lectura} \\
\midrule
C2 & $\log\mu_{ij}=\lambda_{\rho(i)\tau(j)}$ & Un nivel por bloque \\
\midrule
C5 & $\log\mu_{ij}=\alpha_i+\beta_j+\lambda_{\rho(i)\tau(j)}$ & Un nivel por bloque, más un efecto de fila y uno de columna \\
\midrule
C6 & $\log\mu_{ij}=\alpha_{i,\tau(j)}+\beta_{\rho(i),j}$ & El efecto de cada fila cambia según el grupo de columnas, y el de cada columna según el grupo de filas \\
\bottomrule
\end{tabular}
\end{center}
Un modelo \textbf{log-lineal} suma efectos en el logaritmo, y al deshacer el
logaritmo los efectos se \textbf{multiplican}. Por eso C5 es ``media del bloque
$\times$ factor de fila $\times$ factor de columna'': es exactamente la fórmula
reagrupada de la 2.4.3. Aquí se ve de dónde sale esa forma.

\subsection{Las aproximaciones son los estimadores de máxima verosimilitud}

En un modelo log-lineal de Poisson, las medias ajustadas deben reproducir los
totales observados de cada término del modelo. Esos totales son justamente los
promedios que conserva cada base. Entonces las fórmulas de la 2.5.1
\textbf{son} las estimaciones de máxima verosimilitud de esos modelos. Eso cierra
el círculo con el principio MBI de la 2.4.2: la matriz más simple que respeta los
promedios es la que maximiza la verosimilitud de Poisson. El resultado de Birch
garantiza que existe un único conjunto de valores ajustados que cumple el modelo y
reproduce esos totales.

\subsection{Modelos anidados}

Para una partición fija, $\mathrm{C2}\subset\mathrm{C5}\subset\mathrm{C6}$: cada
modelo contiene al anterior como caso particular. C2 es C5 con $\alpha=\beta=0$, y
C5 es C6 reescribiendo los efectos de modo que no dependan del grupo. Es la misma
cadena de la 2.4.3 (de menos a más detalle), ahora vista como modelos
estadísticos.

\subsection{Parámetros libres}

Con $m=n=4$ y $K_R=K_T=2$ (la matriz del ejemplo):
\begin{center}
\begin{tabular}{lcc}
\toprule
Modelo & Fórmula & Valor \\
\midrule
C2 & $K_RK_T$ & 4 \\
C5 & $K_RK_T+(m-K_R)+(n-K_T)$ & 8 \\
C6 & $K_Tm+K_Rn-K_RK_T$ & 12 \\
Saturado & $mn$ & 16 \\
\bottomrule
\end{tabular}
\end{center}
A más parámetros, más flexible el modelo y menor el error, pero más riesgo de
ajustarse al ruido. La diferencia entre el saturado y C6 es
$\sum(|I_g|-1)(|J_h|-1)=4$, igual que los grados de libertad de las pruebas de
independencia por bloque ($16-12=4$).

\begin{idea}
Con I-divergencia, cada base es un modelo de Poisson: C2 un nivel por bloque, C5
añade volumen de fila y columna, C6 permite perfiles propios por bloque. Los tres
están anidados, y $G^2$ mide la asociación que queda dentro de cada bloque.
\end{idea}

\section{2.5.3 Descomposición de la divergencia}

\subsection{Por qué existe esta subsección}

En la 2.5.2 se vio que C2, C5 y C6 son tres modelos anidados, de menos a más
detalle. Cada uno deja un error distinto al aproximar la matriz. La pregunta que
responde esta subsección es: cuando una partición deja un error grande, ¿de qué
está hecho ese error? Es la parte central de la sección, porque justifica por qué
el método de la tesis mide la heterogeneidad con $G^2$ y no con un criterio más
simple.

\subsection{La descomposición}

\[
  D_I(O\|\hat O^{\mathrm{C2}})=
  \underbrace{D_I(\hat O^{\mathrm{C5}}\|\hat O^{\mathrm{C2}})}_{\text{(i) efectos marginales}}+
  \underbrace{D_I(\hat O^{\mathrm{C6}}\|\hat O^{\mathrm{C5}})}_{\text{(ii) marginales por bloque}}+
  \underbrace{\tfrac12\sum_{g,h}G^2_{gh}}_{\text{(iii) interacción local}}.
\]

\textbf{De dónde sale la idea.} Es una propiedad tipo Pitágoras. Si hay dos
modelos anidados, uno simple y uno complejo, ajustados por máxima verosimilitud,
el error del simple se parte en dos:
\begin{center}
error del simple $=$ error del complejo $+$ distancia entre los dos modelos.
\end{center}
Es la misma estructura de la 2.3.2: error con $s$ = error con la media + castigo
por alejarse. Aplicada dos veces a la cadena $\mathrm{C2}\subset\mathrm{C5}\subset\mathrm{C6}$,
y sabiendo que el error de C6 es $\tfrac12\sum G^2$, sale la fórmula.

\subsection{Qué mide cada término}

\begin{center}
\renewcommand{\arraystretch}{1.3}
\begin{tabular}{@{}>{\raggedright\arraybackslash}p{2.6cm}>{\raggedright\arraybackslash}p{6.8cm}>{\raggedright\arraybackslash}p{4.4cm}@{}}
\toprule
\textbf{Término} & \textbf{Qué mide} & \textbf{Pregunta} \\
\midrule
(i) $D(\hat O^{\mathrm{C5}}\|\hat O^{\mathrm{C2}})$ & Cuánto del error se explica solo porque unas regiones e intervalos tienen más volumen que otros & ¿Es heterogeneidad real o solo diferencia de tamaño? \\
\midrule
(ii) $D(\hat O^{\mathrm{C6}}\|\hat O^{\mathrm{C5}})$ & Cuánto cambia el peso relativo de una región según el grupo de horarios (o de un intervalo según el grupo de regiones) & ¿Una región se comporta distinto en bloques distintos? \\
\midrule
(iii) $\tfrac12\sum G^2_{gh}$ & La asociación que queda dentro de cada bloque cuando ya se conocen sus propias marginales & ¿Hay un patrón local dentro del bloque? \\
\bottomrule
\end{tabular}
\end{center}
Ejemplo para (i): una zona de Midtown genera más viajes que una del norte de
Manhattan \emph{a cualquier hora}; eso es volumen, no interacción espacio-temporal.
Ejemplo para (ii): una región que pesa más en los horarios de oficina que en los
nocturnos, aunque su volumen total sea el mismo.

\subsection{Dos ejemplos que aíslan los términos}

\textbf{Ejemplo A: todo es volumen (término i).} La matriz de un solo bloque
$\begin{pmatrix}2&4\\6&12\end{pmatrix}$:
\begin{itemize}
  \item con C2 (todas las celdas valen 6) el error es $\mathbf{4.50}$;
  \item C5 la reproduce exacta, así que el error de C5 es 0;
  \item con un solo bloque, C6 $=$ C5, así que (ii) $=0$ y $G^2=0$;
  \item descomposición: \textbf{(i) $=4.50$, (ii) $=0$, (iii) $=0$.}
\end{itemize}
La matriz parece ``heterogénea'' para un criterio basado en C2, pero no hay ningún
patrón: es un único perfil multiplicado por el volumen de cada fila y columna.

\textbf{Ejemplo B: todo es interacción local (término iii).} La matriz 4$\times$4
de la 2.4:
\begin{itemize}
  \item C5 y C6 coinciden con C2 (las filas de un mismo grupo tienen el mismo
  total), por lo que (i) $=0$ y (ii) $=0$;
  \item $G^2$ por bloque: $1.61$, $2.09$, $0$ y $2.02$, y $\tfrac12\sum G^2=2.86$;
  \item descomposición: \textbf{(i) $=0$, (ii) $=0$, (iii) $=2.86$}, que es justo
  el error total de la 2.4.2.
\end{itemize}
Aquí todo el error es asociación dentro de los bloques: los viajes se concentran
en la diagonal.

\subsection{Por qué esto justifica $G^2$}

\begin{itemize}
  \item Un criterio basado en la constancia del bloque (el implícito en C2) mezcla
  los tres términos. En datos de movilidad con volúmenes muy desiguales, casi todos
  los bloques parecerían heterogéneos solo por el término (i): es el Ejemplo A.
  \item $G^2$ aísla el término (iii), que es lo que se busca: patrones locales de
  asociación.
\end{itemize}
Un bloque de Midtown no es interesante por ser grande; lo es si dentro de él
ciertas regiones se activan en ciertos horarios de forma distinta al resto.

\subsection{La advertencia sobre la prueba $\chi^2$}

$G^2$ tiene, bajo ciertas condiciones, una distribución conocida, la $\chi^2$ con
$(|I_g|-1)(|J_h|-1)$ grados de libertad, que permitiría calcular un valor-p. Por
ejemplo, el bloque $\begin{pmatrix}6&2\\2&6\end{pmatrix}$ tiene $G^2=4.19$ con 1
grado de libertad, lo que daría $p\approx0.04$ y parecería ``significativo''. Pero
esa lectura falla por tres razones:
\begin{enumerate}
  \item \textbf{Conteos pequeños.} La aproximación exige conteos esperados grandes
  y es poco fiable en bloques con poca actividad.
  \item \textbf{La partición se estima con los mismos datos que se prueban.} Se
  eligieron los bloques para que parezcan estructurados y después se pregunta si
  lo son. Eso infla la evidencia (es como elegir el mejor resultado de muchos
  intentos y calcular su significancia como si fuera el único).
  \item \textbf{Sobredispersión.} Los conteos reales de viajes varían más de lo que
  predice Poisson.
\end{enumerate}
Por eso la tesis usa otra forma de calibrar las pruebas.

\begin{idea}
El error de una partición C2 se descompone en volumen (i), cambios de peso por
bloque (ii) e interacción local (iii). $G^2$ aísla el tercero, que es el patrón que
busca la tesis.
\end{idea}

\section{2.5.4 Escala de los datos: conteos frente a promedios}

\subsection{Dos hechos}

\begin{itemize}
  \item \textbf{El co-agrupamiento no depende de la escala.} Por la homogeneidad
  de la I-divergencia (2.3.1), multiplicar $O$ por una constante $c>0$ multiplica
  el error por $c$ y no cambia la mejor partición.
  \item \textbf{La interpretación estadística sí depende de ella.} La equivalencia
  con la devianza de Poisson solo vale si las celdas son \textbf{conteos}.
\end{itemize}

\subsection{Por qué importa: promedios diarios}

Si la matriz contiene promedios diarios, $O=S/D$, donde $S$ son los conteos
totales de $D$ días, entonces $G^2(O)=G^2(S)/D$. Comparar eso con una $\chi^2$
subestima la evidencia por un factor $D$.

\textbf{Con números.} El bloque $\begin{pmatrix}6&2\\2&6\end{pmatrix}$ con
$G^2=4.19$ podría venir de promedios sobre $D=30$ días. Los conteos totales serían
$\begin{pmatrix}180&60\\60&180\end{pmatrix}$ y su $G^2$ sería
$30\times4.19\approx126$. Con 4.19 la evidencia parece débil ($p\approx0.04$); con
126 es abrumadora. Es el mismo patrón, pero la escala cambia por completo la
lectura.

Por eso los resultados de Liu et al.\ (2021), que construyen la matriz con
volúmenes promedio, no permiten una lectura estadística directa. La tesis
construye la matriz con \textbf{conteos totales}, o hace explícito el factor de
escala antes de cualquier prueba.

\begin{idea}
Para que la lectura estadística sea válida, la matriz debe contener conteos
totales, no promedios.
\end{idea}

% =====================================================================
\part{Sección 2.6: El marco ST-CC}

\section{Por qué existe en el capítulo 2}

ST-CC es el marco de Liu et al.\ (2021). \textbf{No es el método de la tesis}: es
el punto de partida que el método propuesto (AST-CC, Capítulo 3) refina. Para
entender qué aporta AST-CC hay que saber qué hace ST-CC y dónde falla. Sus
limitaciones son la motivación de la tesis.

\begin{center}
\renewcommand{\arraystretch}{1.3}
\begin{tabular}{@{}l>{\raggedright\arraybackslash}p{5.4cm}>{\raggedright\arraybackslash}p{4.6cm}@{}}
\toprule
 & \textbf{ST-CC} & \textbf{AST-CC (la tesis)} \\
\midrule
Qué es & El marco de Liu et al.\ (2021) & El método propuesto \\
\midrule
Qué hace & Co-agrupa en tablero con BBAC y reagrupa los bloques con K-means & Parte del tablero de ST-CC y refina cada bloque con pruebas de hipótesis \\
\midrule
Dónde está & Sección 2.6 (marco teórico) & Capítulo 3 \\
\bottomrule
\end{tabular}
\end{center}

\section{Idea general de ST-CC}

Dos etapas, con una matriz región $\times$ intervalo de tiempo:
\begin{enumerate}
  \item \textbf{Co-agrupar en tablero} (con BBAC, 2.4): agrupa filas y columnas a la
  vez y forma bloques.
  \item \textbf{Reagrupar los bloques con K-means}: junta bloques parecidos aunque
  no sean vecinos.
\end{enumerate}
El aporte de los autores no es un algoritmo nuevo, sino la forma de combinar
métodos existentes para interpretar los datos.

\section{2.6.1 Datos y matriz}

\begin{itemize}
  \item \textbf{Datos:} viajes de taxis de Nueva York (noviembre de 2017, unos diez
  millones), por lugar de bajada.
  \item \textbf{Filas:} 66 zonas de Manhattan. \textbf{Columnas:} 12 intervalos de
  dos horas.
  \item \textbf{Dos matrices de $66\times12$:} días laborables (12 días) y
  feriados (9 días).
  \item \textbf{Cada celda:} el \emph{volumen promedio} de bajadas de esa zona en
  ese intervalo.
\end{itemize}
El punto delicado, visto en la 2.5.4: usar promedios no altera la partición, pero
impide la lectura estadística como conteos de Poisson. Por eso AST-CC usa conteos
totales.

\section{2.6.2 Primera etapa: co-agrupamiento en tablero}

\begin{center}
\renewcommand{\arraystretch}{1.3}
\begin{tabular}{@{}>{\raggedright\arraybackslash}p{7.4cm}>{\raggedright\arraybackslash}p{6.2cm}@{}}
\toprule
\textbf{Decisión} & \textbf{Valor} \\
\midrule
Grupos de tiempo $K_T$ & 4 (probaron 4 a 8; el algoritmo devolvía siempre 4) \\
\midrule
Grupos de regiones $K_R$ & 8 (laborables) y 6 (feriados), donde la curva del objetivo se estabiliza \\
\midrule
Inicializaciones & 200 por ejecución \\
\midrule
Convergencia & umbral $10^{-6}$, máximo 2000 iteraciones \\
\bottomrule
\end{tabular}
\end{center}
Resultado: tableros de $8\times4$ y $6\times4$ bloques. En días laborables, los
cuatro grupos de tiempo son madrugada, baja actividad, horario diurno y noche.
Aunque el algoritmo no usa la ubicación, los grupos de regiones salen
espacialmente compactos.

Los valores de $K$ se eligen \textbf{mirando una curva}, sin criterio estadístico.
Esa es una de las limitaciones que ataca la tesis.

\section{2.6.3 Segunda etapa: reagrupar los bloques}

\textbf{El problema:} al asignar filas y columnas completas, bloques distintos
tienen valores parecidos (casi toda la ciudad está tranquila de madrugada).

\textbf{La solución:} cada bloque se describe con \textbf{dos números}, la media y
la desviación estándar de sus celdas (la media resume el nivel de actividad y la
desviación, cuán dispersas están las celdas), y los 32 bloques (laborables) o 24
(feriados) se agrupan con K-means. El número de grupos se elige entre 3 y 8 con
silueta y codo, y sale $K=4$ en ambos periodos. Los cuatro grupos se leen como
niveles de actividad: muy bajo, bajo, medio y alto.

\textbf{Consecuencia:} un grupo reúne bloques de distintas filas y columnas, así
que el resultado deja de ser una rejilla regular. Los autores lo llaman
co-agrupamiento \emph{no en tablero}.

\section{2.6.4 Lectura de los resultados}

\begin{itemize}
  \item \textbf{Patrón espacial:} cómo se distribuyen los niveles de actividad
  sobre todas las regiones dentro de un grupo de horarios. Son a lo sumo cuatro por
  periodo, y su sucesión describe la dinámica del día.
  \item \textbf{Patrón temporal:} la secuencia de niveles que sigue un conjunto de
  regiones a lo largo del día. Encuentran 6 en laborables y 5 en feriados.
  \item \textbf{Hallazgos:} los días laborables son más dinámicos y rítmicos, los
  feriados más estáticos y aleatorios, y Penn Station solo alcanza actividad alta
  en feriados.
\end{itemize}

\section{2.6.5 Limitaciones: las que motivan la tesis}

\begin{center}
\renewcommand{\arraystretch}{1.3}
\begin{tabular}{@{}>{\raggedright\arraybackslash}p{6.6cm}>{\raggedright\arraybackslash}p{6.4cm}@{}}
\toprule
\textbf{Limitación de ST-CC} & \textbf{Qué hace AST-CC} \\
\midrule
K-means fusiona bloques parecidos pero nunca divide uno; la heterogeneidad dentro de un bloque no se recupera & Refinamiento jerárquico: divide los bloques que lo necesitan \\
\midrule
Los números de grupos se eligen con criterios heurísticos (curva, silueta, codo), sin interpretación estadística & Pruebas de hipótesis \\
\midrule
Comparar bloques con media y desviación supone distribuciones semejantes, algo que habría que verificar & Modelo nulo local de cada bloque \\
\bottomrule
\end{tabular}
\end{center}

\section{2.6.6 Base de co-agrupamiento: ST-CC y esta tesis}

Esta subsección resuelve una ambigüedad en Liu et al.:
\begin{center}
\renewcommand{\arraystretch}{1.3}
\begin{tabular}{@{}>{\raggedright\arraybackslash}p{7.2cm}c@{}}
\toprule
\textbf{Qué dicen} & \textbf{Base} \\
\midrule
El nombre del algoritmo (BBAC con I-divergencia) & C2 \\
\midrule
La función objetivo que declaran (pérdida de información mutua) & C5 \\
\midrule
Lo que realmente hace su paquete de código & C2 \\
\bottomrule
\end{tabular}
\end{center}
Entonces ST-CC declara C5 pero implementa C2. Las dos dan un tablero de bloques,
pero no son equivalentes. Por la descomposición de la 2.5.3, C2 mezcla el término
(i) de volumen, así que tiende a agrupar regiones \textbf{por volumen}, mientras
que C5 lo excluye y agrupa \textbf{por forma de perfil}.

La tesis adopta \textbf{C5}, porque el volumen no es interacción espacio-temporal
y porque es lo que los autores declaran. C2 se reporta como variante en los
experimentos. A partir de aquí la subsección deja de describir a Liu et al.\ y
recoge una decisión de la tesis. Sigue pendiente confirmarla con el asesor.

\begin{idea}
ST-CC es el punto de partida: co-agrupa en tablero y reagrupa bloques con K-means.
Sus tres limitaciones son exactamente lo que AST-CC resuelve en el Capítulo 3.
\end{idea}

% =====================================================================
\part{Sección 2.7: Pruebas de hipótesis para tablas de conteo}

\section{Por qué existe esta sección}

En la 2.5 se vio que $G^2$ mide la interacción que queda dentro de un bloque. El
método de la tesis (AST-CC) lo usa para decidir si un bloque se divide. Para eso
hay que resolver tres problemas, y las secciones 2.7 y 2.8 son el kit de
herramientas.

\begin{center}
\renewcommand{\arraystretch}{1.3}
\begin{tabular}{@{}>{\raggedright\arraybackslash}p{9.0cm}>{\raggedright\arraybackslash}p{4.6cm}@{}}
\toprule
\textbf{Problema} & \textbf{Subsección} \\
\midrule
¿Cómo obtener una distribución de referencia \emph{válida} (tablas pequeñas, sobredispersión, grupos elegidos con los mismos datos)? & 2.7.1 a 2.7.3, 2.7.5 \\
\midrule
¿Cómo controlar el error al hacer \emph{muchas pruebas} (una por bloque del árbol)? & 2.7.4 \\
\midrule
¿Cómo evitar que con millones de viajes \emph{cualquier diferencia} resulte significativa? & 2.7.6 \\
\bottomrule
\end{tabular}
\end{center}
La Sección 2.8 presenta SigTree, el método del que la propuesta toma la idea
central: dividir con pruebas de hipótesis y detenerse cuando ninguna división es
significativa. Estas notas cubren 2.7.1 a 2.7.6.

\section{2.7.1 Pruebas de independencia}

\subsection{Los conteos esperados}

La hipótesis nula es que filas y columnas son independientes. Entonces el conteo
esperado de cada celda es
\[
  \hat E_{ij}=\frac{x_{i+}\,x_{+j}}{x_{++}},
\]
es decir, total de la fila por total de la columna entre el total general. Es la
misma fórmula del modelo de independencia de la 2.4.

\subsection{Dos estadísticos}

\begin{itemize}
  \item \textbf{Pearson:} $X^2=\sum_{i,j}\dfrac{(x_{ij}-\hat E_{ij})^2}{\hat E_{ij}}$.
  \item \textbf{Razón de verosimilitud:}
  $G^2=2\sum_{i,j}x_{ij}\log\dfrac{x_{ij}}{\hat E_{ij}}$.
\end{itemize}
Bajo la hipótesis nula y con conteos esperados grandes, ambos siguen aproximadamente
una $\chi^2$ con $(a-1)(b-1)$ grados de libertad. La regla práctica es que la
mayoría de los esperados sean al menos 5.

\textbf{Con el bloque $\begin{pmatrix}6&2\\2&6\end{pmatrix}$:} todos los esperados
valen 4 (menos de 5, ya es una señal de alerta). $X^2=4/4+4/4+4/4+4/4=4$ y
$G^2=4.19$ (calculado en la 2.5). Son parecidos, como se espera. Con 1 grado de
libertad dan $p\approx0.046$ y $p\approx0.041$.

\subsection{Por qué la tesis prefiere $G^2$}

\begin{enumerate}
  \item \textbf{Conecta con Bregman:} $G^2=2D_I(x\|\hat E)$ y
  $G^2=2N\hat I(R;T)$, la información mutua escalada por la masa. Es el mismo
  objeto de la 2.5.
  \item \textbf{Se parte de forma exacta.} Si se agrupan filas, el $G^2$ de la
  tabla completa es igual al $G^2$ de la tabla colapsada más los $G^2$ dentro de
  cada grupo. Pearson no cumple esto.
\end{enumerate}

\textbf{Ejemplo de la propiedad 2.} La tabla con filas $(8,2)$, $(8,2)$ y $(2,8)$
tiene $G^2=10.36$. Se agrupan las dos primeras filas, que son idénticas:
\begin{itemize}
  \item tabla colapsada $\begin{pmatrix}16&4\\2&8\end{pmatrix}$: $G^2=10.36$;
  \item dentro del grupo $\{r_1,r_2\}$: filas idénticas, $G^2=0$; la fila $r_3$
  sola: $G^2=0$;
  \item suma: $10.36+0+0=10.36$. \checkmark
\end{itemize}
Todo lo asociado se explica con ese corte, y no queda nada dentro de los hijos.
Esta es la base de la ``ganancia de un corte'' del Capítulo 3.

\section{2.7.2 Pruebas condicionales y p-valores por simulación}

\subsection{El problema de la $\chi^2$}

La $\chi^2$ es una aproximación y falla en tablas pequeñas, que son frecuentes en
bloques de poca actividad. Hace falta una prueba cuya validez no dependa del
tamaño.

\subsection{La idea: condicionar sobre los totales}

Bajo independencia, si se fijan los totales de fila y de columna, la tabla se
reparte al azar de una forma conocida (la \textbf{hipergeométrica multivariada}: la
forma de repartir los conteos al azar manteniendo esos totales). Esa distribución
no depende de parámetros desconocidos. El p-valor exacto es la probabilidad de
obtener un estadístico al menos tan extremo como el observado. En tablas
$2\times2$ es la prueba exacta de Fisher.

\textbf{Con el mismo bloque $\begin{pmatrix}6&2\\2&6\end{pmatrix}$} (todos los
totales valen 8, $N=16$): las tablas igual o más extremas son las que tienen 0, 1,
2, 6, 7 u 8 en la primera celda. Su probabilidad total es
$1698/12\,870\approx0.13$.
\begin{center}
\begin{tabular}{lc}
\toprule
Método & p-valor \\
\midrule
$\chi^2$ (aproximación) & $\approx0.04$ \\
Fisher (exacto) & $\approx\mathbf{0.13}$ \\
\bottomrule
\end{tabular}
\end{center}
La aproximación decía ``significativo'' y la prueba exacta dice que no hay
evidencia al 5\,\%. Esta es la razón concreta por la que el plan insiste en no
fiarse de la $\chi^2$ en bloques chicos.

\subsection{Simulación (Patefield)}

Enumerar todas las tablas es inviable salvo en tablas muy pequeñas, así que se
simulan $B$ tablas con los mismos totales y se calcula
\[
  \hat p=\frac{1+\#\{k:\,t_k\ge t_{\mathrm{obs}}\}}{1+B}.
\]
El ``+1'' no es un arreglo arbitrario (Phipson y Smyth): la tabla observada es,
bajo la nula, una realización más, y con la corrección la prueba es exacta.
Consecuencias:
\begin{itemize}
  \item el menor p-valor posible es $1/(B+1)$;
  \item para probar a nivel $\alpha$ hace falta $B\ge1/\alpha-1$: $B\ge19$ para
  $\alpha=0.05$ y $B\ge99$ para $\alpha=0.01$;
  \item \textbf{ejemplo:} con $B=99$ y 4 tablas simuladas por encima del observado,
  $\hat p=(1+4)/100=0.05$.
\end{itemize}

\subsection{Simulación secuencial (Besag y Clifford)}

Si hay muchas pruebas y casi todas tienen p grande, simular $B$ tablas en cada una
desperdicia cómputo. Se simula hasta que $h$ tablas superen al observado y se
detiene. Si se detiene tras $l$ tablas, el p-valor es $h/l$. Mantiene la validez y
para pronto en los casos claramente no significativos.

\subsection{El estadístico puede ser cualquiera}

El estadístico puede ser cualquier función de la tabla, incluida una que haga una
búsqueda, por ejemplo el máximo de $G^2$ sobre varias divisiones candidatas. El
p-valor sigue siendo válido siempre que en cada tabla simulada \textbf{se repita la
misma búsqueda}. Es la base del ajuste por selección de la 2.7.5 y del corte
``min-P'' del Capítulo 3.

\section{2.7.3 Remuestreo de réplicas independientes}

\subsection{Por qué hace falta otro método}

La prueba condicional es exacta solo si los conteos son Poisson. Con sobredispersión
no lo son. Hace falta una distribución de referencia que refleje la variabilidad
real, sin suponer un modelo.

\subsection{El bootstrap}

La tabla observada es la suma de $D$ réplicas independientes (los días):
$x=\sum_d x_d$. El bootstrap remuestrea $D$ días \textbf{con reemplazo} y recalcula
el estadístico sobre su suma. Repetido muchas veces, aproxima la distribución del
estadístico.

\textbf{Ejemplo:} con 12 días laborables, cada réplica del bootstrap elige 12 días
al azar con reemplazo (algunos se repiten, otros no salen) y suma esas matrices
diarias.

\subsection{Tres detalles}

\begin{enumerate}
  \item \textbf{Debe hacerse bajo la hipótesis nula.} Aunque los datos no la
  cumplan, se desplazan o reescalan las réplicas para que su media sea la del nulo.
  Si no, la referencia incluiría el efecto que se quiere detectar y la prueba
  perdería potencia.
  \item \textbf{Corrección con pocas réplicas.} La varianza empírica de $D$ réplicas
  subestima la real por el factor $(D-1)/D$, así que las desviaciones se multiplican
  por $\sqrt{D/(D-1)}$.
  \item \textbf{Supuesto que no se puede eludir:} las réplicas deben ser
  \textbf{independientes}. Dos días consecutivos con el mismo episodio de clima no
  lo son, y el remuestreo subestima la variabilidad.
\end{enumerate}

\begin{idea}
$G^2$ se parte de forma exacta y conecta con Bregman, por eso se usa. La $\chi^2$
falla en bloques pequeños ($p\approx0.04$ frente a $0.13$ exacto), así que se usan
pruebas por simulación. Y como los conteos reales están sobredispersos, la
referencia se obtiene remuestreando días.
\end{idea}

\section{2.7.4 Comparaciones múltiples}

\subsection{El problema}

El método hace \textbf{una prueba por cada bloque} del árbol. Si cada una se hace al
nivel $\alpha=0.05$, la probabilidad de equivocarse al menos una vez crece rápido
con el número de pruebas $M$:
\[
  P(\text{al menos un falso rechazo})=1-(1-\alpha)^M.
\]
\begin{center}
\begin{tabular}{cc}
\toprule
$M$ & Probabilidad \\
\midrule
1 & 0.05 \\
20 & 0.64 \\
100 & 0.994 \\
\bottomrule
\end{tabular}
\end{center}
Con 20 pruebas ya es más probable equivocarse que no, aunque no haya ningún efecto
real.

\subsection{Dos formas de medir el error}

\begin{center}
\renewcommand{\arraystretch}{1.3}
\begin{tabular}{@{}l>{\raggedright\arraybackslash}p{6.4cm}>{\raggedright\arraybackslash}p{4.2cm}@{}}
\toprule
\textbf{Medida} & \textbf{Qué controla} & \textbf{Ejemplo} \\
\midrule
FWER & La probabilidad de cometer \emph{al menos un} falso rechazo & Estricta \\
\midrule
FDR & La \emph{proporción esperada} de falsos rechazos entre los rechazos & Con 100 rechazos y FDR 5\,\%, se esperan unos 5 falsos \\
\bottomrule
\end{tabular}
\end{center}
La FWER es más estricta. La FDR tolera algunos falsos a cambio de más potencia, y
suele ser preferible cuando se exploran muchas hipótesis.

\subsection{Las correcciones}

\begin{itemize}
  \item \textbf{Bonferroni:} prueba cada hipótesis al nivel $\alpha/M$. Con $M=20$
  y $\alpha=0.05$, cada prueba usa $0.0025$. Controla la FWER sin supuestos, pero
  es \textbf{muy conservadora} cuando las pruebas están correlacionadas: las
  divisiones candidatas de un mismo bloque comparten casi todos sus datos, así que
  no son 20 pruebas independientes.
  \item \textbf{Min-P (Westfall y Young):} en vez de dividir $\alpha$ entre $M$,
  compara el menor p-valor (o el mayor estadístico) con su distribución nula
  \emph{conjunta}, obtenida por simulación. Tiene en cuenta a la vez el número de
  pruebas y su correlación. Es la que retoma el Capítulo 3.
  \item \textbf{Benjamini y Hochberg:} controla la FDR ordenando los p-valores.
  Vale cuando las pruebas son independientes.
\end{itemize}

\subsection{Pruebas jerárquicas}

Cuando las hipótesis forman un árbol (la de un nodo solo tiene sentido si la del
padre se rechazó), no hace falta corregir por todas las pruebas posibles, sino solo
por las que se hacen. Tres antecedentes:
\begin{itemize}
  \item \textbf{Meinshausen (2008):} de arriba abajo, con un nivel por grupo
  proporcional a su tamaño. Un hijo con el 60\,\% de las variables de su padre
  recibe el 60\,\% del nivel. Controla la FWER sobre toda la jerarquía.
  \item \textbf{Kimes (2017):} lo aplica al agrupamiento jerárquico, recorriendo el
  dendrograma desde la raíz y probando en cada nodo con simulación.
  \item \textbf{Yekutieli (2008):} para FDR, aplica Benjamini y Hochberg dentro de
  cada familia de hijos de una hipótesis rechazada.
\end{itemize}

\subsection{La estructura que usa este problema}

Si la independencia es cierta en un bloque, es cierta en todos sus sub-bloques:
cualquier subtabla de una tabla con medias factorizadas también las tiene
factorizadas. Entonces las hipótesis nulas verdaderas forman \textbf{subárboles
completos}, y un falso rechazo solo puede ocurrir \textbf{por primera vez en la
raíz} de uno de ellos. Por eso basta con controlar ese primer error y no el total
de pruebas.

\section{2.7.5 Inferencia después de la selección}

\subsection{El problema del doble uso de los datos}

Las garantías anteriores suponen que las hipótesis se fijan \textbf{antes} de ver
los datos. En agrupamiento no es así: los grupos, y las hipótesis sobre ellos, se
eligen con los mismos datos.

Gao et al.\ (2024) mostraron que, si los grupos se definen agrupando, una prueba
clásica de diferencia de medias tiene un error de tipo I \textbf{extremadamente
inflado}, y que el problema persiste aunque se usen dos conjuntos de datos
independientes. La razón es que el agrupamiento \textbf{encuentra grupos distintos
incluso cuando no hay estructura}. Si se generan datos de puro ruido, se agrupan en
dos y se prueba si los grupos difieren, la prueba casi siempre sale significativa.

\subsection{Tres remedios}

\begin{center}
\renewcommand{\arraystretch}{1.3}
\begin{tabular}{@{}>{\raggedright\arraybackslash}p{3.8cm}>{\raggedright\arraybackslash}p{5.4cm}>{\raggedright\arraybackslash}p{4.4cm}@{}}
\toprule
\textbf{Enfoque} & \textbf{Cómo funciona} & \textbf{Límite} \\
\midrule
1. Incluir la selección en la simulación & En cada réplica nula se repite el procedimiento de selección completo & El p-valor vale para la nula \emph{global}, no para hipótesis más específicas \\
\midrule
2. Inferencia selectiva & Se calcula el p-valor condicionando al evento de que el procedimiento produjo esos grupos (Gao 2024, Watanabe 2021) & Garantías formales, pero casi siempre para datos gaussianos y costoso \\
\midrule
3. División de cada observación (count splitting) & Ver abajo & Falla con sobredispersión \\
\bottomrule
\end{tabular}
\end{center}

\subsection{Count splitting (Neufeld et al., 2024)}

Cada conteo se divide en dos partes, $x=x^{(1)}+x^{(2)}$, con
$x^{(1)}\mid x\sim\mathrm{Binomial}(x,\varepsilon)$. Si $x\sim\mathrm{Poisson}(\mu)$,
las dos partes son \textbf{independientes} y Poisson, con medias $\varepsilon\mu$ y
$(1-\varepsilon)\mu$. La estructura se elige con $x^{(1)}$ y se prueba con
$x^{(2)}$, como si se hubiera fijado de antemano.

\textbf{Ejemplo:} una celda con $\mu=10$ y $\varepsilon=0.5$ produce
$x^{(1)}\sim\mathrm{Poisson}(5)$ y $x^{(2)}\sim\mathrm{Poisson}(5)$, independientes.
Se elige el corte con la primera mitad y se prueba con la segunda.

Es lo que distingue este enfoque de dividir las \emph{filas} en dos grupos de datos,
que no funciona según Gao: aquí cada conteo se parte, no cada fila. Pero si los
conteos están sobredispersos, las dos partes dejan de ser independientes y la
garantía se debilita. \emph{Data thinning} generaliza la idea a otras familias.

\section{2.7.6 Tamaño de efecto y sobredispersión}

\subsection{Tamaño de efecto: el problema de los millones de viajes}

Con muestras grandes, una prueba detecta desviaciones mínimas. Como
$G^2=2N\hat I(R;T)$, si la información mutua poblacional es positiva, aunque sea
ínfima, $G^2$ crece \textbf{proporcional a $N$} y el p-valor tiende a 0.

\textbf{Ejemplo:} con $\hat I=0.001$ nats, $G^2=2$ con $N=1\,000$, pero
$G^2=2\,000$ con $N=1\,000\,000$. Misma asociación diminuta; con un millón de viajes
es ``altamente significativa''. Los autores de SigTree reconocen este fenómeno como
limitación.

La solución es acompañar la significancia con una medida de tamaño de efecto que no
crezca con $N$:
\begin{center}
\renewcommand{\arraystretch}{1.3}
\begin{tabular}{@{}>{\raggedright\arraybackslash}p{4.4cm}>{\raggedright\arraybackslash}p{5.2cm}>{\raggedright\arraybackslash}p{3.6cm}@{}}
\toprule
\textbf{Medida} & \textbf{Fórmula} & \textbf{Rango} \\
\midrule
Información mutua por observación & $\hat I=G^2/(2N)$ (nats) & Conecta con la función objetivo \\
\midrule
$V$ de Cramér & $\sqrt{X^2/(N\min(a-1,b-1))}$ & 0 (independencia) a 1 (asociación perfecta) \\
\bottomrule
\end{tabular}
\end{center}

\textbf{Con números.} El bloque $\begin{pmatrix}6&2\\2&6\end{pmatrix}$: $X^2=4$,
$N=16$, $\min=1$, así que $V=\sqrt{4/16}=0.5$ e $\hat I=4.19/32=0.131$. Multiplicado
por 100, $\begin{pmatrix}600&200\\200&600\end{pmatrix}$: $G^2=419$ (100 veces
mayor), pero $V=0.5$ e $\hat I=0.131$ \textbf{iguales}. La significancia explota con
$N$; el tamaño de efecto no.

\subsection{El sesgo de la información mutua empírica}

Aunque la población sea independiente, $G^2$ sigue una $\chi^2$ con
$\nu=(a-1)(b-1)$ grados de libertad, cuya media es $\nu$. Entonces
$\mathrm{E}[\hat I]\approx\nu/(2N)>0$. Es el sesgo que corrige Miller y Madow para
estimadores de información mutua.
\begin{center}
\begin{tabular}{cccc}
\toprule
Tabla & $\nu$ & $N$ & Sesgo esperado de $\hat I$ \\
\midrule
$2\times2$ & 1 & 100 & 0.005 \\
$5\times5$ & 16 & 100 & 0.08 \\
\bottomrule
\end{tabular}
\end{center}
Un umbral sobre $\hat I$ sin corregir favorecería a las tablas con pocas
observaciones y muchas celdas, porque tendrían un $\hat I$ ``alto'' solo por azar.

\subsection{Sobredispersión}

Poisson impone que la varianza sea igual a la media. En datos reales es mayor,
$\mathrm{Var}=\varphi\mu$ con $\varphi>1$, porque los eventos no son
independientes: el día de la semana, el clima, eventos puntuales. Entonces $X^2$ y
$G^2$ se \textbf{inflan por el factor $\varphi$} y la prueba rechaza demasiado.

\textbf{Corrección por cuasi-verosimilitud} (ajusta la varianza sin suponer una
distribución completa): dividir el estadístico entre una estimación de $\varphi$
(por ejemplo, el cociente de Pearson entre sus grados de libertad).

\textbf{Ejemplo:} con $\varphi=4$, un estadístico de 8 con 1 grado de libertad
($p\approx0.005$) se corrige a 2 ($p\approx0.16$). Pasa de ``muy significativo'' a
``no significativo''.

\subsection{Funciones de varianza}

Un único $\varphi$ supone varianza proporcional a la media. Hay formas más generales
$\mathrm{Var}=V(\mu)$:
\begin{itemize}
  \item \textbf{Binomial negativa:} $V(\mu)=\mu+\mu^2/k$. La sobredispersión crece
  con la media. Con $k=50$: $\mu=10$ da varianza 12 ($\varphi=1.2$), pero $\mu=100$
  da 300 ($\varphi=3$).
  \item \textbf{Ley de potencia de Taylor:} $V(\mu)=a\mu^b$, que contiene como casos
  particulares Poisson ($a=b=1$) y $\varphi$ constante ($b=1$).
\end{itemize}

\subsection{El límite de todos estos modelos}

Todas describen la varianza de \textbf{cada celda por separado}. Cuando la
variabilidad viene de factores que actúan sobre muchas celdas a la vez (el clima de
un día), las celdas dejan de ser independientes entre sí, y ningún modelo de
varianza por celda lo captura. Si hay réplicas independientes (los días), la
correlación se reproduce \textbf{remuestreando} (el bootstrap de la 2.7.3) en lugar
de modelarla.

\begin{idea}
Con muchas pruebas el error se acumula (64\,\% con 20), y como los bloques de un
árbol están anidados basta con controlar el primer rechazo falso. Elegir grupos con
los mismos datos infla el error y se corrige repitiendo la selección en la
simulación o partiendo cada conteo. Con millones de viajes todo es significativo,
así que se exige además un tamaño de efecto corregido por sesgo. La sobredispersión
infla los estadísticos y se corrige remuestreando días.
\end{idea}

% =====================================================================
\part{Sección 2.8: SigTree}

\section{Por qué existe esta sección}

SigTree (Hu et al., 2025) es el método del que la propuesta toma su idea central:
\textbf{construir una partición con pruebas de hipótesis y detenerse cuando ninguna
división es significativa}. Importa por una razón concreta. En ST-CC (2.6) el
número de grupos se elige mirando una curva, sin criterio estadístico. SigTree no
necesita fijar el número de grupos: tiene un solo parámetro, el nivel $\alpha$, con
interpretación estadística clara. Eso es justo lo que busca la tesis.

La sección tiene dos funciones: explica cómo funciona SigTree y deja claro qué hay
que adaptar para trasladarlo al problema de la tesis (bloques de una matriz de
conteos en vez de muestras con atributos categóricos). Las herramientas de la 2.7
son las piezas con las que se hace esa adaptación.

\subsection{Antecedente: CHAID}

Usar pruebas de hipótesis para decidir divisiones de un árbol no es nuevo.
\textbf{CHAID} (Kass, 1980) lo hace en el caso \textbf{supervisado}: elige cada
división con una prueba $\chi^2$ entre la división candidata y la variable de
clase. SigTree traslada la idea al caso \textbf{no supervisado}, donde no hay
variable de clase.

\section{2.8.1 Divisiones candidatas y pseudo-clases}

\textbf{El escenario.} $n$ muestras descritas por $m$ atributos categóricos (el
atributo $q$ con $k_q$ categorías).

\textbf{La división candidata.} En cada nodo, cada categoría $c$ de cada atributo
$j$ define una división: separa las muestras con valor $c$ en el atributo $j$ de las
que no. Se representa con una variable binaria $Z=\mathbb 1[x_j=c]$.

\textbf{El problema:} no hay variable de clase contra la cual probar la división.
\textbf{La solución:} cada uno de los otros $m-1$ atributos hace de
\textbf{pseudo-clase}, es decir, de etiqueta provisional.
\begin{itemize}
  \item \textbf{Hipótesis nula de cada prueba:} $Z$ es independiente del atributo
  $q$.
  \item \textbf{Hipótesis nula global:} $Z$ es independiente de todos los demás
  atributos.
  \item \textbf{Intuición:} si las muestras del nodo no tienen estructura de
  grupos, ninguna división candidata estará asociada con la mayoría de las
  pseudo-clases. Si sí hay grupos, una buena división separará a la vez muchos
  atributos.
\end{itemize}

\section{2.8.2 Estadístico y criterio de división}

Para cada par (división $Z$, atributo $q\ne j$):
\begin{enumerate}
  \item se construye la tabla de contingencia $2\times k_q$;
  \item se calcula el estadístico de Pearson $X^2(Z,q)$; en tablas $2\times2$ con
  algún esperado menor que 5 se aplica la \textbf{corrección de Yates}, un ajuste
  que evita sobrestimar el estadístico con conteos pequeños;
  \item se suman todos los atributos:
  \[
    X^2(Z)=\sum_{q\ne j}X^2(Z,q),\qquad \mathrm{gl}(Z)=\sum_{q\ne j}(k_q-1);
  \]
  \item el p-valor se obtiene de la $\chi^2$ con $\mathrm{gl}(Z)$ grados de
  libertad, suponiendo que las pruebas individuales son independientes;
  \item en cada nodo se elige la división de \textbf{menor p-valor}.
\end{enumerate}

\subsection{Ejemplo numérico}

Un nodo con 100 muestras y tres atributos binarios $A$, $B$ y $C$. Se prueba la
división $Z=\mathbb 1[A=1]$: 50 muestras tienen $A=1$ y 50 no.
\begin{center}
\renewcommand{\arraystretch}{1.3}
\begin{tabular}{@{}lccc@{}}
\toprule
Pseudo-clase & Tabla $Z\times$ atributo & Esperados & $X^2$ \\
\midrule
$B$ & $\begin{pmatrix}40&10\\10&40\end{pmatrix}$ & 25 en cada celda & $4\cdot15^2/25=\mathbf{36}$ \\
\midrule
$C$ & $\begin{pmatrix}30&20\\20&30\end{pmatrix}$ & 25 en cada celda & $4\cdot5^2/25=\mathbf{4}$ \\
\bottomrule
\end{tabular}
\end{center}
La suma es $X^2(Z)=36+4=40$ con $\mathrm{gl}=1+1=2$. El p-valor de la $\chi^2$ con 2
grados de libertad es $e^{-20}\approx2\cdot10^{-9}$, muy pequeño. La división
$A=1$ está fuertemente asociada con $B$ y débilmente con $C$, y la suma lo recoge.
SigTree compara este p-valor con el de las demás candidatas y se queda con la menor.

\section{2.8.3 Criterios de parada}

La división elegida solo se aplica si su p-valor es menor o igual que un nivel
ajustado
\[
  \alpha^\ast=\frac{\alpha}{|Q|},
\]
donde $|Q|$ es el número total de categorías de todos los atributos. Es una
corrección de Bonferroni con un divisor \textbf{constante}, igual en todos los nodos
(2.7.4). En el ejemplo, con tres atributos binarios, $|Q|=6$ y
$\alpha^\ast=0.05/6\approx0.0083$. Como $2\cdot10^{-9}<0.0083$, la división se
acepta.

Además hay dos reglas triviales: no se dividen nodos con 5 muestras o menos, y no
se aceptan divisiones que produzcan un hijo con menos de 3 muestras. El crecimiento
termina cuando ninguna hoja admite una división significativa, y las hojas son los
grupos finales. El resultado es un algoritmo con \textbf{un solo parámetro}
interpretable, el nivel $\alpha$, que \textbf{no requiere fijar el número de
grupos}.

\section{2.8.4 Limitaciones}

\textbf{Las tres que reconocen los autores:}
\begin{enumerate}
  \item Evaluar todas las divisiones candidatas en cada nodo es costoso.
  \item No optimiza una función objetivo explícita, así que no hay garantía de que
  el árbol sea el mejor posible.
  \item El estadístico $\chi^2$ crece con el número de muestras y de atributos, lo
  que empuja los p-valores a cero y dificulta controlar los falsos descubrimientos
  con ese umbral (el problema de la 2.7.6).
\end{enumerate}
\textbf{Dos que se añaden desde la 2.7:}
\begin{enumerate}
  \setcounter{enumi}{3}
  \item El p-valor de la división elegida se calcula \emph{como si esa división se
  hubiera fijado de antemano}, sin tener en cuenta que se eligió por ser la de menor
  p-valor (el problema de la 2.7.5).
  \item La suma de estadísticos supone independencia entre pruebas que
  \emph{comparten la misma variable} $Z$, lo que no es cierto.
\end{enumerate}

\section{2.8.5 Relación con la propuesta}

La tabla compara componente por componente. Leída con lo visto en la 2.7, cada fila
apunta a una herramienta:
\begin{center}
\renewcommand{\arraystretch}{1.3}
\footnotesize
\begin{tabular}{@{}>{\raggedright\arraybackslash}p{2.7cm}>{\raggedright\arraybackslash}p{3.3cm}>{\raggedright\arraybackslash}p{4.6cm}>{\raggedright\arraybackslash}p{2.6cm}@{}}
\toprule
\textbf{Componente} & \textbf{SigTree} & \textbf{Propuesta} & \textbf{Herramienta (2.7)} \\
\midrule
Objeto que se divide & Muestras con atributos categóricos & Bloque de una matriz de conteos región $\times$ tiempo & (adaptación al problema) \\
\midrule
División candidata & Presencia o ausencia de una categoría & Bipartición de las filas, de las columnas o de ambas & (adaptación) \\
\midrule
Qué se prueba & Asociación entre la división y los demás atributos & Asociación entre la división y la \emph{otra dimensión} del bloque & 2.7.1 \\
\midrule
Estadístico & Suma de $X^2$ & $G^2$ de la tabla colapsada, coherente con la I-divergencia & 2.7.1 ($G^2$ se parte de forma exacta) \\
\midrule
Distribución de referencia & $\chi^2$ asintótica & \emph{Bootstrap} de los días alrededor de la independencia del bloque, repitiendo la búsqueda & 2.7.3 y 2.7.5 \\
\midrule
Comparaciones múltiples & Bonferroni uniforme & Reparto jerárquico del nivel a lo largo del árbol & 2.7.4 \\
\midrule
Criterio de parada & Significancia y reglas triviales & Significancia, tamaño de efecto corregido por sesgo y reglas triviales & 2.7.6 \\
\bottomrule
\end{tabular}
\end{center}
Las limitaciones de ambos métodos orientan el diseño. La elección heurística de
$K_R$ y $K_T$ en ST-CC motiva que la resolución final no dependa solo de esos
números, y la inflación del estadístico de SigTree con muestras grandes motiva
complementar la significancia con un umbral de tamaño de efecto.

\begin{idea}
SigTree divide un árbol con pruebas de hipótesis y para cuando nada es
significativo, sin fijar el número de grupos. La propuesta conserva ese principio y
cambia cada pieza: $G^2$ en lugar de la suma de $X^2$, \emph{bootstrap} de días en
lugar de $\chi^2$, reparto jerárquico del nivel en lugar de Bonferroni uniforme, y
tamaño de efecto además de significancia.
\end{idea}

% =====================================================================
\part{Capítulo 3, bloque 1: visión general y formulación}

\section{Por qué existe el capítulo}

El Capítulo 2 dejó un problema abierto. ST-CC resume cada bloque de su tablero con
un solo valor, así que la interacción entre regiones y horarios que queda
\emph{dentro} de los bloques, el término (iii) de la descomposición de la
sección 2.5, no se modela. Un patrón local, por ejemplo unas pocas regiones que se
activan solo en la tarde, queda diluido en un bloque que parece uniforme. Para
recuperarlo hacen falta dos cosas, una forma de subdividir un bloque sin tocar el
resto de la matriz y un criterio estadístico para decidir si la subdivisión es
real y cuándo parar. El Capítulo 2 aportó las herramientas ($G^2$, pruebas por
simulación, control jerárquico, tamaño de efecto) y el Capítulo 3 las junta.

\section{3.1 Visión general}

La idea cabe en una frase. Se refina mientras haya heterogeneidad local detectable
y de tamaño relevante, y se para cuando todas las hojas sean indistinguibles de su
modelo nulo. Es el principio de SigTree llevado a una matriz. El método tiene
cuatro etapas, que son construir la matriz $O$ con sus $D$ matrices diarias,
obtener la partición base con BBAC$_I$ y la base C5, refinar cada bloque
(dividir por filas, por columnas, de forma conjunta, o dejarlo como hoja) y
resumir las hojas con su tamaño de efecto y su nivel de actividad.

La diferencia con subir $K_R$ y $K_T$ es que un corte global atraviesa toda la
matriz y fragmenta también las zonas sin patrón, mientras que un corte de AST-CC
afecta solo al bloque donde se detectó la heterogeneidad. Los parámetros crecen
con la estructura local encontrada y no con la resolución global.

\section{3.2 Formulación}

\subsection{Notación y la cantidad central}

Un bloque es $B=I\times J$, con $I$ un subconjunto de regiones y $J$ uno de
horarios. Se calcula su tabla esperada bajo independencia,
$E^B_{ij}=O^B_{i\cdot}O^B_{\cdot j}/N_B$, y su estadístico
$G^2_B=2\sum O_{ij}\log(O_{ij}/E^B_{ij})$, que es la sección 2.5 aplicada a un
bloque. La información mutua empírica es $\hat I_B=G^2_B/(2N_B)$.

\paragraph{Ejemplo.} Sea el bloque de cuatro regiones y dos horarios
\[
\begin{array}{c|cc}
 & \text{mañana} & \text{tarde}\\\hline
A & 30 & 10\\ B & 28 & 12\\ C & 10 & 30\\ D & 12 & 28
\end{array}
\]
Cada fila suma 40 y cada columna 80, con $N_B=160$, así que todos los esperados
valen 20. Entonces
$G^2_B=2\sum O\ln(O/20)\approx 34.10$ con $(4-1)(2-1)=3$ grados de libertad, y
$\hat I_B=34.10/320\approx 0.107$ nats por viaje. Se ve que $A$ y $B$ se parecen,
que $C$ y $D$ también, y que los dos pares son opuestos.

\subsection{Corte, árbol y modelo local}

Un corte parte un bloque de tres maneras, de filas (dos hijos), de columnas (dos
hijos) o conjunto (cuatro hijos). El árbol de refinamiento tiene como hijos de la
raíz los $K_RK_T$ bloques base, y debajo los cortes. Las hojas forman siempre una
partición de la matriz con estructura de guillotina, porque cada corte actúa solo
dentro de su bloque, y si no se corta nada coinciden con el tablero de ST-CC.

Cada hoja se aproxima con el modelo de independencia de su propio bloque,
$\hat O=E^L$. La ventaja es que cada hoja tiene sus propios efectos de fila y de
columna, así que las diferencias de volumen entre regiones no cuentan como patrón.
El número de parámetros es $\sum_L(|I_L|+|J_L|-1)$.

\subsection{La proposición clave}

Tiene tres partes. En (a), la divergencia del árbol es la mitad de la suma de las
$G^2$ de sus hojas. En (b), si se corta $B$ con $s$, entonces
$G^2_B=G^2_{B,s}+\sum_{C}G^2_C$, un reparto exacto entre lo que explica el corte y
lo que queda dentro de los hijos. En (c), cortar reduce la divergencia en
exactamente $\tfrac12G^2_{B,s}$, y $G^2_{B,s}$ es un estadístico de razón de
verosimilitud con grados de libertad conocidos ($|J|-1$ para filas, $|I|-1$ para
columnas, $|I|+|J|-3$ para conjunto).

\paragraph{Ejemplo con el bloque anterior.} Se cortan las filas en $\{A,B\}$ y
$\{C,D\}$. La tabla colapsada es
\[
\begin{array}{c|cc}
 & \text{mañana} & \text{tarde}\\\hline
\{A,B\} & 58 & 22\\ \{C,D\} & 22 & 58
\end{array}
\]
con esperados 40. Con la fórmula de la proposición,
$\log\frac{O_{rj}N_B}{N_rO^B_{\cdot j}}=\log\frac{58\cdot160}{80\cdot80}=\ln1.45$,
y
\[
G^2_{B,s}=2\cdot2\,[58\ln1.45+22\ln0.55]\approx 33.59
\]
con $|J|-1=1$ grado de libertad. Cada hijo es casi homogéneo, por ejemplo
$(30,10;\,28,12)$ con esperados $(29,11)$ y $G^2_C\approx0.25$; los dos suman
$0.51$. Se comprueba el reparto, $33.59+0.51\approx34.10=G^2_B$. El corte explica el
98.5\,\% de la heterogeneidad ($33.59/34.10$) con un solo parámetro más, pues
antes hay $4+2-1=5$ parámetros y después $3+3=6$.

\paragraph{Por qué importan los grados de libertad.} Un corte de columnas separa
las dos columnas y explica el 100\,\%, pero con $|I|-1=3$ grados de libertad.
Con $G^2$ crudo parecería mejor, y por eso la comparación entre tipos de corte se
hace con p-valores calibrados.

\paragraph{Lectura informacional.} Dividiendo el reparto entre $2N_B$ se obtiene la
regla de la cadena de la información mutua, así que la asociación del bloque se
reparte sin pérdida ni duplicación entre lo que explica el corte y lo que queda en
los hijos.

\paragraph{Por qué funciona (no va en la tesis).} Dentro de cada hoja, la
aproximación reproduce el total, así que los términos lineales de la I-divergencia
se cancelan y queda $\sum O\log(O/E)=\tfrac12G^2_L$. Para el reparto se usa la
misma identidad pitagórica de la sección 2.5, aplicada a dos modelos anidados, el
de independencia en $B$ y el del corte. Su diferencia es la ganancia
$G^2_{B,s}$, que resulta ser el $G^2$ de la tabla colapsada porque, tras un corte de
filas, los totales de fila de cada hijo coinciden con los de $B$.

\subsection{Relación con ST-CC (el corolario)}

ST-CC deja sin modelar el término (iii), la suma de las $G^2$ de los bloques base.
AST-CC lo reparte en una parte \emph{explicada} (las ganancias de los cortes
aceptados) y una parte \emph{residual} (lo que queda en las hojas). Cada corte
aceptado traslada $\tfrac12G^2_{B,s}$ de la segunda a la primera. El reparto
depende solo de la partición y no de la base con que se obtuvo, así que vale con C2
y con C5, y el criterio de refinamiento y el objetivo de ST-CC son la misma
cantidad a distintas escalas.

\subsection{Planteamiento del problema}

El problema es minimizar la divergencia del árbol sujeto a que ninguna hoja tenga
heterogeneidad significativa y relevante. No se resuelve de forma exacta porque el
árbol se construye de forma voraz, de arriba abajo. La proposición garantiza que
cada decisión local es coherente con el objetivo global, no que el árbol final sea
óptimo. La versión penalizada tipo BIC, con $\lambda=\tfrac12\log N$, no construye
el árbol sino que sirve para comparar en una misma escala árboles de AST-CC con
rejillas globales de distinto tamaño.

\begin{idea}
Cortar un bloque reparte su $G^2$ de forma exacta entre lo que explica el corte y lo
que queda en los hijos. Por eso cada corte se puede decidir con la información de
su propio bloque sin perder coherencia con el objetivo global, y su ganancia es un
estadístico de razón de verosimilitud al que se le puede aplicar una prueba de
hipótesis.
\end{idea}

% =====================================================================
\part{Capítulo 3, bloque 2: estadísticos y regla de decisión}

\section{Por qué existen estas secciones}

La sección 3.2 mostró que cortar un bloque reparte su $G^2$ de forma exacta, pero
no dijo cuándo conviene cortar ni dónde. La 3.3 define las preguntas estadísticas
que se le hacen a cada bloque, y la 3.4 las junta en una regla que termina en una de
tres decisiones, que son cortar, dejar el bloque como hoja homogénea o dejarlo como
hoja heterogénea no localizada.

\section{3.3 Estadísticos de heterogeneidad}

\subsection{3.3.1 Elección del modelo nulo}

Hay dos candidatos para lo que se considera comportamiento esperado dentro de un
bloque. La \emph{constancia} supone que todas las celdas valen lo mismo,
$\mu_{ij}=\theta$, y es el nulo implícito en la base C2. La \emph{independencia}
supone $\mu_{ij}=a_ib_j$, un efecto de región multiplicado por un efecto de horario.

\paragraph{Ejemplo.} Sea el bloque con filas $(80,20)$ y $(40,10)$, con $N=150$. La
segunda región tiene la mitad del volumen de la primera pero el mismo perfil, pues
ambas ponen el 80\,\% en la mañana. No hay ningún patrón local.
\begin{itemize}
\item Bajo constancia el esperado es $37.5$ en todas las celdas y
$G^2=2\,[80\ln\tfrac{80}{37.5}+20\ln\tfrac{20}{37.5}+40\ln\tfrac{40}{37.5}+10\ln\tfrac{10}{37.5}]\approx74.8$
con 3 grados de libertad, que se rechaza con claridad.
\item Bajo independencia los esperados son exactamente $(80,20)$ y $(40,10)$, así
que $G^2=0$.
\end{itemize}
Con datos de movilidad, donde las zonas tienen volúmenes muy distintos, la
constancia rechazaría casi todos los bloques solo por el volumen, que son los
términos (i) y (ii) de la descomposición. La independencia aísla el término (iii),
y rechazarla significa que algunas regiones se activan en algunos horarios de forma
distinta al resto, que es la definición operativa de patrón local de esta tesis.

\subsection{3.3.2 Prueba ómnibus}

La primera pregunta es si el bloque contiene algo que refinar, y se responde con
$G^2_B$. Con conteos de Poisson independientes bastaría una simulación condicionada
a los totales, pero los taxis tienen una varianza entre días mucho mayor que la
media. Con var/media $\approx3$, el valor típico de $G^2_B$ bajo el nulo ronda
$3\nu$ en lugar de $\nu$. Con $\nu=12$ serían unas 36 unidades en lugar de 12, y una
prueba $\chi^2$ con 12 grados de libertad daría $p<0.001$ en un bloque que no tiene
nada. Por eso el nulo se genera remuestreando los días reales (sección 3.5).

\subsection{3.3.3 Pruebas dirigidas}

Si el ómnibus rechaza, la segunda pregunta es qué corte captura la estructura. Para
cada tipo se toma el máximo de la ganancia $G^2_{B,s}$ sobre los candidatos. Dentro
de un tipo todos los candidatos tienen los mismos grados de libertad, así que el de
mayor $G^2$ es el de menor p-valor, igual que SigTree elige la división de menor
p-valor.

\paragraph{La observación del nulo común.} El bloque es independiente si y solo si
todas sus filas tienen el mismo perfil sobre las columnas, de modo que el ómnibus y
las dirigidas contrastan la misma hipótesis con estadísticos de distinta potencia.
Las dirigidas solo se hacen si el ómnibus rechazó, así que para que un bloque nulo
se refine hay que rechazar primero el ómnibus, lo que ocurre con probabilidad a lo
sumo $\alpha_B$. Las dirigidas localizan la estructura y no añaden oportunidades de
falso rechazo.

\paragraph{Un caso que parece problemático.} Sean las filas $(30,10)$, $(10,30)$,
$(30,10)$ y $(10,30)$. Todos los totales de fila valen 40 y los de columna 80, así
que las marginales no muestran nada. Sin embargo el corte $\{1,3\}$ frente a
$\{2,4\}$ colapsa el bloque en $(60,20)$ y $(20,60)$, con esperados 40, y
\[
G^2_{B,s}=2\cdot2\,[60\ln1.5+20\ln0.5]\approx41.9 .
\]
El corte compara los perfiles completos de los grupos y no sus totales. Por eso los
cortes unidimensionales bastan para la mayoría de las estructuras, y el conjunto se
reserva para cuando la evidencia está repartida entre ambas dimensiones.

\subsection{3.3.4 Generación de cortes candidatos}

Un bloque con $|I|$ filas admite $2^{|I|-1}-1$ biparticiones. Con 4 filas son 7,
con 14 son 8191 y con 15 serían 16383, así que cada fila adicional duplica el
número de candidatos.

\paragraph{Columnas.} Con $n=12$ hay a lo sumo $2^{11}-1=2047$ biparticiones, y se
enumeran todas.

\paragraph{Filas.} Si $|I|\le14$ se enumeran. Por encima se usa un 2-medias con
divergencia KL ponderada. Maximizar la ganancia de un corte de filas equivale a
minimizar $\sum_iO^B_{i\cdot}\,\mathrm{KL}(p(\cdot|i)\,\|\,p(\cdot|r(i)))$, donde
$p(\cdot|i)$ es el perfil horario de la fila $i$. Es un K-means de dos grupos que
agrupa las filas por la forma de su perfil y no por su volumen, y es el paso de filas
del co-agrupamiento de Dhillon con dos grupos. Cada fila va al grupo cuyo perfil
promedio se parece más al suyo, y se repite hasta que nadie cambia, desde
$R_{\mathrm{ini}}=10$ inicios aleatorios.

\paragraph{Conjunto.} La ganancia de un corte conjunto equivale a minimizar la suma
de las $G^2$ de los cuatro hijos. Si $|J|\le6$ y $|I|\le14$ se recorren todos los
cortes de columnas y, para cada uno, se elige el mejor de filas, lo que da el óptimo
exacto. En bloques mayores se alternan ambos pasos.

\paragraph{Una búsqueda heurística no invalida la prueba.} Siempre que el mismo
procedimiento se repita en cada réplica nula, el p-valor es válido para el estadístico
que efectivamente se calcula, aunque no sea el máximo exacto. Lo que se pierde es
potencia, y por eso se prefiere enumerar siempre que se pueda.

\section{Por qué el límite es 14 filas}

El 14 no es un número teórico sino un compromiso entre exactitud y costo, y
conviene tener claras sus tres piezas.

\paragraph{1. Por qué enumerar.} El 2-medias es una heurística que puede quedarse en
un mínimo local. En una comprobación sobre 40 tablas de $8\times5$, en las que el
óptimo se calcula enumerando, lo alcanzó en alrededor del 92\,\% de los casos con
diez inicios, y nunca lo superó, porque una heurística no puede ganarle al óptimo.
Ese 8\,\% de fallos es lo que justifica enumerar siempre que se pueda. Importa
porque el p-valor compara el máximo observado con el máximo bajo el nulo, y un
máximo mal encontrado resta potencia.

\paragraph{2. Por qué no enumerar siempre.} Cada fila adicional duplica los
candidatos, y la búsqueda no se hace una vez sino en cada réplica nula. Con
$|I|=14$ son 8191 candidatos, cada uno una tabla colapsada $2\times|J|$. Con las
$B_{\mathrm{perm}}$ réplicas, que van de 199 a 20\,000 según el nivel del bloque, y
las tres búsquedas (filas, columnas y conjunta), el trabajo ya es considerable. Con
15 filas se duplicaría, y con 20 filas serían más de medio millón de candidatos por
réplica. Además la implementación procesa las réplicas por lotes en memoria, y el
lote que cabe se reduce a la mitad por cada fila adicional.

\paragraph{3. Por qué 14 y no 13 o 15.} El valor concreto es una elección práctica,
no un umbral que salga de una fórmula. Es el tamaño hasta el cual la enumeración
sigue siendo manejable con el número de réplicas de referencia. Lo que sí se
comprobó es que, con este límite, todos los bloques base de los datos de esta tesis,
que tienen a lo sumo quince filas, se resuelven de forma casi siempre exacta, pues
solo los de quince filas, si los hay, usan la heurística.

\paragraph{Advertencia honesta.} El 92\,\% se midió con tablas de ocho filas, no
de quince. Con más filas el 2-medias podría fallar algo más, y eso es lo que
justifica no extender la heurística a bloques grandes. Si el jurado pregunta, la
respuesta es que 14 es un compromiso de costo, que el fallo de la heurística se
midió en tablas pequeñas, y que la validez de la prueba no depende de que la búsqueda
sea exacta, solo de que se repita igual en las réplicas.

\section{3.4 Regla de decisión}

La regla combina el ómnibus, las dirigidas y el umbral de tamaño de efecto $\delta$,
con $\alpha_B$ el nivel asignado al bloque.
\begin{enumerate}
\item \textbf{Puerta ómnibus.} Si $p_{\mathrm{omn}}>\alpha_B$ o $\tilde I_B<\delta$,
el bloque es una hoja homogénea.
\item \textbf{Cortes unidimensionales.} Si solo uno de $p_{\mathrm{fil}}$,
$p_{\mathrm{col}}$ es $\le\alpha_B$, se aplica el mejor corte de esa dimensión. Si lo
son los dos, se aplica el de menor p-valor y la recursión decide sobre la otra. Los
empates se resuelven con $z=(G^2-\nu)/\sqrt{2\nu}$.
\item \textbf{Corte conjunto.} Si ninguno es significativo pero el ómnibus sí lo fue,
se evalúa el $2\times2$. Si $p_{\mathrm{conj}}\le\alpha_B$ se aplica, y si no, el
bloque es una hoja heterogénea no localizada.
\item \textbf{Ganancia mínima.} El corte elegido solo se acepta si
$\tilde g_{B,s}\ge\delta$. Si no, el bloque también queda como hoja heterogénea no
localizada.
\end{enumerate}

\paragraph{Cuatro trayectorias con números.} Sea $\alpha_B=0.01$ y $\delta=0.005$.
\begin{itemize}
\item \textbf{A.} $p_{\mathrm{omn}}=0.2$. El ómnibus no rechaza y el bloque es hoja
homogénea.
\item \textbf{B.} $p_{\mathrm{omn}}=0.002$, $\tilde I_B=0.03\ge\delta$,
$p_{\mathrm{fil}}=0.004$ y $p_{\mathrm{col}}=0.3$. Solo el corte de filas es
significativo, se aplica, y si $\tilde g\ge\delta$ se recurre en cada hijo.
\item \textbf{C.} Ambos p-valores valen $0.001$, el mínimo con 999 réplicas. El
corte de filas tiene $G^2=60$ con $\nu=5$ y el de columnas $G^2=45$ con $\nu=9$.
Entonces $z_{\mathrm{fil}}=55/\sqrt{10}\approx17.4$ y
$z_{\mathrm{col}}=36/\sqrt{18}\approx8.5$, y gana el corte de filas.
\item \textbf{D.} El ómnibus rechaza pero ni los cortes unidimensionales ni el
conjunto son significativos. El bloque queda como hoja heterogénea no localizada.
\end{itemize}

\paragraph{Las tres decisiones de diseño.}
\begin{itemize}
\item \emph{Un corte por nivel.} Cortar ambas dimensiones a la vez pasaría de uno a
cuatro bloques cuando quizá bastan dos. Cortar una por nivel añade menos parámetros
y da una lectura direccional, pues si el bloque se divide por filas la
heterogeneidad es espacial, y si se divide por columnas, temporal.
\item \emph{El corte conjunto como recurso.} Su modelo contiene a los de filas y de
columnas, así que su estadístico es al menos tan grande como cualquiera de ellos y
puede alcanzar significancia cuando la estructura está repartida entre ambas
dimensiones.
\item \emph{Dos tipos de hoja.} La homogénea no muestra interacción. La heterogénea
no localizada sí la muestra, pero ningún corte simple la captura. Reportarlas por
separado evita presentar como homogéneo un bloque que no lo es.
\end{itemize}

\begin{idea}
La regla es una cascada de puertas. Primero se pregunta si hay algo (ómnibus y
tamaño de efecto), luego dónde (filas, columnas, conjunto) y por último si vale la
pena (ganancia mínima). El error de tipo I lo controla la primera puerta, y las
demás solo localizan. El límite de 14 filas es un compromiso entre exactitud de la
búsqueda y costo, y no afecta la validez de la prueba.
\end{idea}

% =====================================================================
\part{Capítulo 3, bloque 3: calibración de las pruebas}

\section{Por qué existe esta sección}

Todas las pruebas de la 3.3 y la 3.4 terminan pidiendo un p-valor, y un p-valor solo
vale si se sabe cómo se comporta el estadístico cuando no hay patrón. Con conteos
de Poisson independientes esa distribución se conoce, pero los datos de taxis no lo
son. La sección responde a una pregunta práctica, que es cómo fabricar tablas nulas
que se parezcan a los datos reales salvo por la ausencia de interacción. Tiene tres
partes: medir qué tan lejos está Poisson de la realidad, construir la nula con
\emph{bootstrap} de días y corregir que el corte probado fue elegido como el mejor.

\section{3.5.1 Variabilidad entre días}

La matriz es la suma de $D$ días, así que cada celda tiene réplicas diarias
$X_{dij}$ y se puede contrastar Poisson directamente. Con Poisson, var/media valdría
1 en todas las celdas.

\begin{center}
\small
\begin{tabular}{@{}lcc@{}}
\toprule
 & Laborable ($D=12$) & Feriado ($D=9$) \\
\midrule
Mediana de var/media & 3.2 & 17.1 \\
Con media diaria en $[5,50)$ & 2.0 & 2.8 \\
Con media diaria $\ge500$ & 4.3 & 49.1 \\
Exponente $b$ de $\mathrm{var}=a\mu^b$ & 1.22 & 1.81 \\
Dispersión de Pearson tras quitar efectos día--región y día--intervalo & 2.1 & 6.7 \\
\bottomrule
\end{tabular}
\end{center}

\paragraph{Hecho 1: la sobredispersión es sustancial.} Una celda típica tiene
varianza 3 veces su media en día laborable y 17 veces en feriado. Con la nula de
Poisson, la prueba rechazaría la independencia solo por esto.

\paragraph{Hecho 2: no es constante y crece con el volumen.} Un modelo
cuasi-Poisson usa un único factor $\varphi$. Con $\varphi=17.1$, la mediana del
feriado, una celda poco activa con razón $2.8$ se corrige de más unas 6 veces
($17.1/2.8$), y una muy activa con razón $49.1$ se corrige de menos casi 3 veces
($49.1/17.1$). Las muy activas son las que concentran la masa del bloque.

\paragraph{Hecho 3: una parte no afecta a la prueba y otra es correlacionada.} Si un
día multiplica la intensidad de toda una región por un factor, la suma sobre los
días sigue siendo un producto de efecto de fila por efecto de columna.

\emph{Ejemplo.} Sea $a=(1,2)$ para dos regiones y $b=(10,20)$ para dos horarios. El
día 1 tiene $\gamma=(1,1)$ y el día 2 tiene $\gamma=(2,0.5)$, es decir, en el día 2 la
región 1 está muy activa y la 2 poco. Sumando los dos días, la región 1 queda con
$1\cdot(1+2)\cdot(10,20)=(30,60)$ y la región 2 con $2\cdot(1+0.5)\cdot(10,20)=(30,60)$.
Las filas son proporcionales, así que el bloque sigue siendo independiente y no hay
patrón falso.

Lo que sí queda es la combinación de efectos día--región y día--intervalo, que
produce desviaciones correlacionadas entre las celdas de un bloque con forma de
interacción, y cuyo peso en $G^2$ crece con la masa del bloque. Una nula que trate
cada celda por separado no puede reproducirlas, y por eso se remuestrean días
completos.

\section{3.5.2 \emph{Bootstrap} de días}

\paragraph{El p-valor.} Se generan $B$ réplicas nulas y se cuenta cuántas tienen un
estadístico mayor o igual al observado, $p=(1+\#)/(1+B)$. Con $B=199$ y 4 réplicas
por encima, $p=5/200=0.025$. Si ninguna supera, $p=1/200=0.005$. Nunca vale 0, y
el mínimo posible es $1/(B+1)$.

\paragraph{La réplica.} Para cada celda con media diaria $\bar X\ge0.5$ se calcula
la fluctuación de cada día, $F_d=X_d/\bar X$, y luego se corrige,
$\tilde F=\max\{0,\,1+\sqrt{D/(D-1)}\,(F-1)\}$. La corrección compensa que la
varianza empírica de $D$ días subestima la poblacional.

\emph{Ejemplo.} Una celda con $D=4$ días y conteos $(80,100,90,130)$, con media 100.
Las fluctuaciones son $F=(0.8,1.0,0.9,1.3)$. El factor es $\sqrt{4/3}=1.1547$ y las
corregidas quedan $\tilde F=(0.769,\,1.0,\,0.885,\,1.346)$. Las desviaciones al
cuadrado suman $0.14$, que dividido entre $D$ da $0.035$ y dividido entre $D-1$ da
$0.0467$. La corrección multiplica la varianza por $4/3$ y la deja insesgada.

\emph{Una réplica.} Se sortean 4 días con reemplazo, por ejemplo los días
$(4,1,4,2)$. La suma de sus $\tilde F$ es $1.346+0.769+1.346+1.0=4.461$. Si la
esperada de independencia del bloque es $E^B=400$, la réplica vale
$400/4\cdot4.461=446.1$, y se redondea de forma estocástica a 446 o 447 sin sesgo.

\paragraph{Por qué es una nula válida.} La media de $\tilde F$ sobre los días es
exactamente 1, así que la media de las réplicas es $E^B$ y se cumple el nulo. Si un
día solo difiere por un factor común a una región o a un intervalo, su fluctuación
es de rango uno y la réplica sigue siendo independiente. La combinación de efectos y
la variación propia de cada celda se trasladan tal como aparecen en los datos. El
método no supone ninguna forma para la varianza ni para la correlación, pero sí que
los días son independientes y que $E^B$ estima bien las medias.

\paragraph{Evidencia de calibración.} Para cada bloque base, 300 réplicas generadas
con la fórmula anterior cumplen la independencia por construcción y conservan la
variabilidad real. Error de tipo~I al nivel 0.05, como media sobre los 24 bloques.

\begin{center}
\small
\begin{tabular}{@{}lcc@{}}
\toprule
Nula & Laborable & Feriado \\
\midrule
Poisson & 0.40 & 0.76 \\
Cuasi-Poisson & 0.03 & 0.01 \\
Gamma--Poisson & 0.02 & 0.06 \\
\emph{Bootstrap} sin corrección & 0.06 & 0.11 \\
\emph{Bootstrap} con corrección & \textbf{0.04} & \textbf{0.09} \\
\bottomrule
\end{tabular}
\end{center}

Poisson rechaza una independencia verdadera en el 40\,\% y el 76\,\% de los casos.
Las nulas por celda son conservadoras en la mayoría de los bloques y liberales en los
de mayor volumen. El \emph{bootstrap} corregido es el único centrado en el nominal en
el periodo laborable.

\paragraph{El feriado.} El generador mezcla al azar sábados, domingos y el Día de
Acción de Gracias, mientras la matriz real tiene siempre la misma composición. Con un
generador que la respeta, el error medio baja a 0.045, aunque 5 de los 24 bloques,
casi todos nocturnos, siguen por encima, con máximo 0.30. Remuestrear dentro de cada
tipo de día empeora a 0.11, porque con solo cuatro sábados y cuatro domingos no hay
datos suficientes. Por eso se usa el \emph{bootstrap} sin estratos y la liberalidad
residual se declara como limitación. Como estos datos nulos comparten mecanismo con
el \emph{bootstrap}, lo que podría favorecerlo, la calibración final se medirá
también con datos sintéticos de otro mecanismo.

\section{3.5.3 Corrección por la búsqueda del mejor corte}

El corte que se prueba se eligió porque maximiza la separación. Si su $G^2$ se
comparara con una $\chi^2$ de un corte fijado de antemano, el p-valor estaría
sesgado hacia cero.

\paragraph{Intuición.} Con 7 cortes candidatos independientes de 1 grado de
libertad, la probabilidad de que al menos uno supere el umbral del 5\,\% de una
$\chi^2$ es $1-0.95^7\approx\mathbf{0.30}$ y no $0.05$. Los candidatos reales están
correlacionados, así que es menos, pero el sesgo sigue ahí, y con las 2047
biparticiones de columnas es mucho mayor. Es el doble uso de los datos que
documentan Gao y colaboradores en el agrupamiento jerárquico.

\paragraph{La corrección tipo \emph{min-P}.} En cada réplica nula se repite la
búsqueda completa del mejor corte y se registra su máximo. El p-valor compara el
máximo observado con la distribución del máximo bajo el nulo. Ajusta a la vez por el
número de candidatos y por su correlación, con más potencia que Bonferroni. En la
comprobación preliminar, comparar el mejor corte con una $\chi^2$ fija dio errores de
tipo~I de hasta 0.98 al nivel 0.05, y con \emph{min-P} el error medio fue 0.050. Las
pruebas dirigidas reutilizan las mismas réplicas del ómnibus, así que el costo de
generarlas se paga una sola vez por bloque.

\section{3.5.4 Alcance de la corrección}

La repetición de la búsqueda corrige la selección del corte dentro de un nodo, pero
no la selección del propio nodo. Los bloques base vienen de BBAC$_I$, que arma los
grupos de modo que sean lo más distintos posible entre sí, y los hijos vienen de
cortes elegidos por separar bien los mismos datos con que luego se prueban. Aunque
cada prueba individual esté bien calibrada, el error del método completo puede
ser mayor que el nominal, y por eso el capítulo de resultados incluye un experimento
sobre datos nulos que mide cuántos refinamientos espurios produce el método completo.

\begin{idea}
El p-valor solo es honesto si la nula se parece a los datos reales. Remuestrear días
completos conserva la correlación entre celdas que Poisson y las nulas por celda
pierden, y repetir la búsqueda del mejor corte en cada réplica evita el sesgo de
haberlo elegido con los mismos datos. Lo que queda sin corregir es la selección de
los propios nodos, y eso solo lo mide el experimento con datos nulos.
\end{idea}

% =====================================================================
\part{Capítulo 3, bloque 4: comparaciones múltiples y criterios de parada}

\section{Por qué existen estas secciones}

Un árbol de refinamiento no hace una prueba sino muchas, una por cada bloque base y,
en cascada, una por cada hijo. Si todas usan $\alpha=0.05$, la probabilidad de algún
falso refinamiento crece rápido. La sección 3.6 reparte el nivel de significancia
para que el árbol completo tenga un error global controlado. La 3.7 resuelve otro
problema, que es que con millones de viajes casi cualquier desviación resulta
significativa, y por eso se exige además un tamaño de efecto mínimo y se fijan reglas
para detenerse.

\section{3.6 Control de comparaciones múltiples}

\paragraph{El problema.} Con 20 pruebas independientes al 5\,\%, la probabilidad de
al menos un falso rechazo es $1-0.95^{20}\approx64$\,\%. SigTree divide $\alpha$
entre el número de pruebas (Bonferroni), pero ignora la estructura jerárquica y
resulta demasiado conservador.

\paragraph{El reparto por masa.} El nivel global se reparte de arriba abajo en
proporción a la masa de cada bloque. Para los bloques base,
$\alpha_B=\alpha N_B/N$, y para los hijos, $\alpha_C=\alpha_BN_C/N_B$.

\emph{Ejemplo.} Sea $N=10\,000$ viajes y $\alpha=0.05$, con cuatro bloques base de
4000, 3000, 2000 y 1000 viajes. Sus niveles son $0.020$, $0.015$, $0.010$ y $0.005$,
que suman $0.05$. Si el bloque de 4000 se corta en hijos de 3000 y 1000, estos
reciben $0.02\cdot0.75=0.015$ y $0.02\cdot0.25=0.005$, que suman lo que tenía el
padre. Con Bonferroni sobre 20 pruebas todos recibirían $0.0025$, y el bloque de 4000
trabaja con $0.02$, ocho veces más, lo cual es razonable porque tiene más datos y más
potencia. Además, un bloque solo se prueba si su padre fue refinado, y los niveles de
los bloques que nunca se visitan no se gastan.

\paragraph{El argumento de control.} Si el nulo de independencia se cumple en un
bloque, se cumple en todos sus sub-bloques, así que los nulos verdaderos forman
subárboles completos. Un falso rechazo ocurre por primera vez en la raíz de alguno de
esos subárboles, y esas raíces son bloques disjuntos cuyos niveles suman a lo sumo
$\alpha$. En el ejemplo, si los bloques de 2000 y 1000 viajes fueran verdaderamente
nulos, la probabilidad de refinarlos por error sería a lo sumo
$0.010+0.005=0.015\le0.05$. Eso da control de la tasa de error por familia (FWER),
la probabilidad de al menos un falso refinamiento en todo el árbol. La salvedad
honesta es que la jerarquía se construye con los datos, así que la garantía es
aproximada, y su comportamiento real lo mide el experimento de calibración del
capítulo de resultados.

\paragraph{Consecuencia práctica, el número de réplicas.} El nivel de un bloque
pequeño puede ser menor que la resolución del p-valor por simulación, $1/(B+1)$. Por
eso $B_{\mathrm{perm}}=\max\{B_{\min},\lceil c/\alpha_B\rceil\}$, acotado por
$B_{\max}$, con $B_{\min}=199$, $c=2$ y $B_{\max}=20\,000$.

\begin{center}
\small
\begin{tabular}{@{}cccl@{}}
\toprule
$\alpha_B$ & $\lceil2/\alpha_B\rceil$ & $B_{\mathrm{perm}}$ & Resultado \\
\midrule
0.02 & 100 & 199 & Se usa el mínimo \\
0.005 & 400 & 400 & Se adapta \\
0.0001 & 20\,000 & 20\,000 & Llega al máximo \\
0.00002 & 100\,000 & 20\,000 & $p_{\min}=5\cdot10^{-5}>\alpha_B$, hoja inalcanzable \\
\bottomrule
\end{tabular}
\end{center}

Una hoja inalcanzable no es una hoja homogénea. El método no afirma que el bloque sea
homogéneo, sino que no puede probarlo al nivel asignado.

\paragraph{Parada secuencial de Besag y Clifford.} La simulación se detiene en cuanto
se acumulan $h=10$ réplicas con estadístico mayor o igual al observado. Si el bloque
claramente no es significativo, por ejemplo con $p\approx0.5$, eso ocurre tras unas 20
réplicas en lugar de 199. Los bloques que sí pasan el ómnibus completan la
simulación, y esas réplicas sirven para las pruebas dirigidas.

\section{3.7 Criterios de parada}

\paragraph{Significancia frente a tamaño de efecto.} Un bloque con $N=10^6$ viajes e
$\hat I=0.0005$ tiene $G^2=2\cdot10^6\cdot0.0005=1000$, un p-valor prácticamente
nulo y una interacción irrelevante. Con $\delta=0.005$ es hoja. Otro con $N=1000$ e
$\hat I=0.05$ tiene $G^2=100$, y si pasa la prueba también supera $\delta$, así que se
refina. La significancia crece con $N$ y el tamaño de efecto no.

\paragraph{Los siete criterios.}

\begin{center}
\small
\begin{tabular}{@{}p{3.2cm}p{4.6cm}p{5.4cm}@{}}
\toprule
Criterio & Condición para ser hoja & Cuándo importa \\
\midrule
Tamaño mínimo & $|I|<2$ o $|J|<2$ & Con una sola fila o columna, $G^2=0$ \\
Masa mínima & $N_B<N_{\min}$ (500) & Bloques con muy poca actividad \\
Profundidad máxima & $d\ge d_{\max}$ (6) & Acota costo y complejidad \\
Nivel inalcanzable & $1/(B_{\max}+1)>\alpha_B$ & La prueba no podría rechazar \\
Significancia & $p_{\mathrm{omn}}>\alpha_B$ o ningún corte $\le\alpha_B$ & No hay evidencia o no se localiza \\
Tamaño de efecto & $\tilde I_B<\delta$ & Interacción irrelevante \\
Ganancia del corte & $\tilde g<\delta$ & El corte explica muy poco \\
\bottomrule
\end{tabular}
\end{center}

Los cuatro primeros son reglas triviales, análogas a las de SigTree, que no dependen
de la estructura interna del bloque. Los tres últimos son estadísticos.

\paragraph{El umbral $\delta$.} Es la defensa principal contra la inflación por tamaño
de muestra y tiene una interpretación directa, información mutua en nats por
observación, sin crecer con $N$. En los bloques base de ST-CC el cociente $G^2/\nu$
va de 3 a más de 1700 y la prueba de independencia rechaza en casi todos. Por eso, en
los primeros niveles del árbol decide $\delta$, y la prueba cobra importancia en los
bloques pequeños y profundos, donde la masa es menor.

\paragraph{Tamaño de efecto corregido por sesgo.} Aun sin interacción,
$\mathrm{E}[G^2]\approx\nu$, así que $\mathrm{E}[\hat I]\approx\nu/(2N)>0$. Un bloque
de 1033 viajes con $\nu=12$ tiene $\hat I=0.017$, y cerca de un tercio,
$12/2066\approx0.006$, es ruido de Poisson. La versión corregida es
$\tilde I=\max\{0,(G^2-\bar t_0)/(2N)\}$, con $\bar t_0$ la media de $G^2$ bajo las
réplicas nulas.

\emph{Ejemplo ilustrativo.} Aquí $G^2=2\cdot1033\cdot0.017\approx35.1$. Restando el
$\nu=12$ de Poisson, $\tilde I=(35.1-12)/2066\approx0.011$. Con sobredispersión de
2.1, como en el periodo laborable, $\bar t_0$ sería más cercano a 25 y entonces
$\tilde I=(35.1-25)/2066\approx0.005$. Los 25 son una cifra supuesta solo para mostrar
el efecto. El mismo bloque pasa de $0.017$ a $0.011$ y a $0.005$ según cuánto ruido se
descuente. La ganancia corregida $\tilde g$ funciona igual, pero con $\bar t_{0,s}$
calculado repitiendo la búsqueda, así que descuenta además el optimismo de haber
elegido el mejor corte. Ambas medias salen de las réplicas que ya se generaron para
los p-valores, sin costo adicional.

\paragraph{Elección de $\delta$.} No se fija a priori. Los $\hat I$ de los bloques
base van de $0.0012$ a $0.061$ en laborable y de $0.0017$ a $0.19$ en feriado. En el
capítulo de resultados se recorre una rejilla como
$\{0.001,0.0025,0.005,0.01,0.02\}$, y para cada valor se reportan el número de hojas,
la fracción del término de interacción local que explica el árbol y la estabilidad de
las hojas. Como medida estandarizada complementaria puede reportarse la $V$ de
Cramér.

\begin{idea}
El nivel $\alpha$ se reparte por masa a lo largo del árbol, de modo que los bloques con
más datos tienen más poder de rechazo y el error global queda acotado. Con muchos
datos casi todo es significativo, por eso la significancia se acompaña de un umbral de
tamaño de efecto $\delta$, corregido por el sesgo que produce el ruido, y de reglas
triviales de parada.
\end{idea}

% =====================================================================
\part{Capítulo 3, bloque 5: el algoritmo}

\section{Por qué existe esta sección}

Hasta ahora el método se presentó por piezas, que son el reparto de $G^2$, la prueba
ómnibus, las pruebas dirigidas, la nula con \emph{bootstrap}, el reparto del nivel y
los criterios de parada. La sección 3.8 las junta en dos pseudocódigos que se pueden
implementar, fija los valores de referencia de los parámetros y estima el costo, para
mostrar que el método es ejecutable en datos reales.

\section{3.8 Algoritmo}

\subsection{Algoritmo 1: Refinar}

La entrada es la matriz $O$ con sus $D$ días y los parámetros. Primero se obtiene la
partición base con BBAC$_I$ y luego se llama a \textsc{Refinar} sobre cada bloque base
con nivel $\alpha N_B/N$ y profundidad 0. \textsc{Refinar} sigue seis pasos.
\begin{enumerate}
\item \textbf{Reglas triviales.} Si $|I|<2$, $|J|<2$, $N_B<N_{\min}$ o
$d\ge d_{\max}$, el bloque es hoja.
\item \textbf{Número de réplicas.} Se calcula $B_{\mathrm{perm}}$. Si
$1/(B_{\mathrm{perm}}+1)>\alpha_B$, es hoja inalcanzable.
\item \textbf{Puerta ómnibus.} Se calcula $G^2_B$, se generan las réplicas por
\emph{bootstrap} de días (con parada de Besag y Clifford) y se obtienen $p_{\mathrm{omn}}$,
$\bar t_0$ y $\tilde I_B$. Si $p_{\mathrm{omn}}>\alpha_B$ o $\tilde I_B<\delta$, es
hoja homogénea.
\item \textbf{Pruebas dirigidas.} Se llama a \textsc{MejorCorte} para filas y para
columnas con las mismas réplicas. Si el menor p-valor es $\le\alpha_B$, se elige ese
corte (desempate por $z$). Si no, se prueba el conjunto, y si tampoco es
significativo, es hoja heterogénea no localizada.
\item \textbf{Ganancia mínima.} Se calcula $\tilde g$. Si $\tilde g<\delta$, es hoja
heterogénea no localizada.
\item \textbf{Recursión.} Se llama a \textsc{Refinar} sobre cada hijo con
$\alpha_C=\alpha_BN_C/N_B$ y profundidad $d+1$.
\end{enumerate}

\paragraph{Seguimiento con números.} Sea $N=10\,000$, $\alpha=0.05$ y un bloque base
de 4000 viajes con $\alpha_B=0.02$ y $\delta=0.005$.
\begin{itemize}
\item Pasa las reglas triviales, pues tiene filas y columnas de sobra, $4000\ge500$ y
$d=0$.
\item $B_{\mathrm{perm}}=\min\{20\,000,\max\{199,\lceil2/0.02\rceil=100\}\}=199$. El
p mínimo es $1/200=0.005\le0.02$, así que es alcanzable.
\item Supón $p_{\mathrm{omn}}=0.005$ y $\tilde I=0.03\ge\delta$. Pasa la puerta.
\item Supón $p_{\mathrm{fil}}=0.005$ y $p_{\mathrm{col}}=0.4$. Solo el corte de filas
es significativo y se elige.
\item Supón $\tilde g\ge\delta$. Se acepta y se recurre en cada hijo. Si el corte dejó
hijos de 3000 y 1000 viajes, reciben $0.015$ y $0.005$. El hijo de nivel $0.005$ usa
$B_{\mathrm{perm}}=\max\{199,400\}=400$ y repite todo el proceso.
\end{itemize}
La recursión termina cuando todas las ramas llegan a una hoja, y el resultado es el
árbol con cada hoja etiquetada como homogénea, heterogénea no localizada o
inalcanzable.

\paragraph{Una precisión sobre las réplicas.} La parada secuencial se activa solo
cuando se acumulan $h=10$ réplicas por encima del observado, lo que implica un
p-valor mayor que $\alpha_B$ (con $B_{\mathrm{perm}}\approx2/\alpha_B$, esas 10
réplicas equivalen a un p-valor de al menos $5\alpha_B$). Por eso, si el bloque pasa la
puerta ómnibus, la simulación ya estaba completa y \textsc{MejorCorte} puede
reutilizarla sin pasos adicionales.

\subsection{Algoritmo 2: \textsc{MejorCorte} con \emph{min-P}}

Recibe un bloque, un tipo de corte y las réplicas nulas, y hace cuatro cosas.
\begin{enumerate}
\item Busca el mejor corte en el bloque observado y calcula
$t_{\mathrm{obs}}=G^2_{B,s^\ast}$.
\item \textbf{Repite la búsqueda completa en cada réplica}, y no solo evalúa el corte
encontrado en los datos reales. Esa repetición es lo que corrige el sesgo de haber
elegido el mejor corte.
\item Calcula $p=(1+\#\{t_b\ge t_{\mathrm{obs}}\})/(1+B_{\mathrm{perm}})$.
\item Devuelve el corte, el p-valor y la media $\bar t$ de los $t_b$, que alimenta la
corrección por sesgo de $\tilde g$.
\end{enumerate}

\subsection{La tabla de parámetros}

Los parámetros son de dos tipos. Las \emph{decisiones del investigador}, que se
someten a análisis de sensibilidad, son $K_R$ y $K_T$ (resolución inicial, tomados de
Liu), $\alpha=0.05$, $\delta$ (rejilla en $[0.001,0.02]$ nats), $N_{\min}=500$ y
$d_{\max}=6$. Los \emph{parámetros técnicos}, que dan fiabilidad a la prueba sin ser el
foco de la interpretación, son las 200 inicializaciones de BBAC$_I$,
$B_{\min}=199$, $c=2$, $B_{\max}=20\,000$, $h=10$, los límites de enumeración
$|I|\le14$ y $|J|\le6$, y $R_{\mathrm{ini}}=10$ del 2-medias. Todos estos valores
coinciden con los de la implementación (\texttt{astcc/arbol.py}).

\paragraph{Una nota de consistencia.} El criterio de nivel inalcanzable de la tabla de
parada usa $B_{\max}$ y el algoritmo usa $B_{\mathrm{perm}}$, que está acotado por
$B_{\max}$. Son equivalentes, porque si $\alpha_B$ es grande la prueba es alcanzable
con $B_{\mathrm{perm}}$ y si es pequeño $B_{\mathrm{perm}}$ llega a $B_{\max}$.

\subsection{Costo computacional}

\paragraph{Partición base.} Cuesta lo mismo que ST-CC, $O(mn(K_R+K_T))$ por iteración.
La búsqueda local evalúa cada movimiento en $O(K_RK_T)$, porque con las bases C2 y C5
la divergencia solo depende de los totales de bloque.

\paragraph{Refinamiento.} El costo extra está en las pruebas por simulación. Por
bloque, el orden es
\[
B_{\mathrm{perm}}\,\bigl[\,2^{|J|}|I|+2^{\min(|I|,14)}|J|\,\bigr],
\]
donde el primer término es la búsqueda de columnas y el segundo la de filas. Con
$n=12$ y bloques base de a lo sumo 15 filas, quedan acotados.

\begin{center}
\small
\begin{tabular}{@{}lccc@{}}
\toprule
Bloque & Términos por réplica & Con 199 réplicas & Con 20\,000 \\
\midrule
$|I|=10$, $|J|=12$ & $2^{12}\cdot10+2^{10}\cdot12\approx5\cdot10^4$ & $\approx10^7$ & $\approx10^9$ \\
$|I|=14$, $|J|=12$ & $2^{12}\cdot14+2^{14}\cdot12\approx2.5\cdot10^5$ & $\approx5\cdot10^7$ & $\approx5\cdot10^9$ \\
\bottomrule
\end{tabular}
\end{center}

Son órdenes de magnitud de operaciones elementales, no tiempos medidos. Los bloques con
20\,000 réplicas son los de nivel muy pequeño y poca masa, que se vuelven hojas pronto.

\paragraph{Tres propiedades mantienen el costo manejable.} La parada secuencial
abandona pronto los bloques no significativos. Los bloques de un mismo nivel son
independientes dado el árbol, por lo que se procesan en paralelo. Y el tamaño de los
bloques decrece con la profundidad.

\begin{idea}
El algoritmo es una recursión que, en cada bloque, hace tres preguntas en orden (si
hay algo, dónde y si vale la pena), reutiliza las mismas réplicas nulas para todas
ellas y reparte el nivel $\alpha$ a los hijos en proporción a su masa. El costo extra
frente a ST-CC está solo en las simulaciones, y la parada secuencial, el
paralelismo por niveles y la reducción de tamaño lo mantienen manejable.
\end{idea}

% =====================================================================
\part{Capítulo 3, bloque 6: representación de los resultados}

\section{Por qué existe esta sección}

ST-CC se lee fácil porque la partición es un tablero, con $K_T$ patrones espaciales
con su dinámica temporal y $K_R$ patrones temporales con su distribución espacial.
Con refinamiento local esa lectura deja de existir como un solo objeto, porque la
partición temporal puede cambiar de un subconjunto de regiones a otro. Si el método
devolviera solo un árbol de hojas sería correcto pero difícil de interpretar. La
sección define cinco productos que devuelven la interpretabilidad que se pierde al
abandonar la rejilla.

\section{3.9 Los cinco productos}

\subsection{1. Descriptores de hoja}

Para cada hoja se reportan sus regiones y horarios, su masa $N_L$, su interacción
residual corregida $\tilde I_L$, su tipo (homogénea, heterogénea no localizada o
inalcanzable) y su nivel de actividad relativo
\[
\ell_L=\frac{N_L}{\sum_{(i,j)\in L}O_{i\cdot}O_{\cdot j}/N}.
\]
Compara la actividad que se observa en la hoja con la que predecirían los volúmenes
globales de sus regiones y de sus horarios. Como la hoja es un rectángulo, el
denominador se simplifica a
$\bigl(\sum_{i\in I_L}O_{i\cdot}\bigr)\bigl(\sum_{j\in J_L}O_{\cdot j}\bigr)/N$.

\paragraph{Ejemplo.} Sea $N=10\,000$ y una hoja con las regiones 1 y 2 y los horarios
$a$ y $b$. Los totales globales son $O_{1\cdot}=500$, $O_{2\cdot}=300$,
$O_{\cdot a}=1000$ y $O_{\cdot b}=2000$. El esperado es
$(500+300)(1000+2000)/10\,000=800\cdot3000/10\,000=\mathbf{240}$. Si la hoja tiene
$N_L=360$ viajes, entonces $\ell_L=360/240=\mathbf{1.5}$. Ese grupo de regiones se
activa en ese grupo de horarios un 50\,\% más de lo que corresponde a su volumen. Un
valor menor que 1 indica lo contrario, que está más apagado de lo esperado.

\paragraph{Por qué la referencia es la matriz completa y no el padre.} El modelo de
independencia del bloque padre reproduce los totales de fila y de columna del padre.
Tras un corte de filas, el esperado de cada hijo suma exactamente su masa, y el
cociente respecto del padre vale 1 en todos los hijos, sin distinguir nada. En el
bloque de 4000 viajes cortado en hijos de 3000 y 1000, respecto del padre ambos
darían $3000/3000=1$ y $1000/1000=1$. Con la matriz completa como referencia, el
cociente sí mide algo.

\subsection{2. Árbol anotado}

Cada nodo interno muestra el tipo de corte, la ganancia corregida $\tilde g$ y el
p-valor, de modo que se lee qué parte de la heterogeneidad fue espacial y cuál
temporal. Cada hoja se lee como el resultado de la secuencia de cortes que llevó hasta
ella, como las reglas de un árbol de decisión.

\paragraph{Ejemplo de lectura.} Un bloque base se corta por filas ($\tilde g=0.04$,
$p=0.005$) y uno de sus hijos se corta por columnas ($\tilde g=0.01$, $p=0.01$). La
hoja resultante se lee como un grupo de regiones, dentro de ese bloque, en un grupo
de horarios. La primera pregunta fue espacial y la segunda temporal.

\subsection{3. Matriz teselada}

La matriz se reordena según la partición base y, dentro de cada bloque, según sus
cortes. Se dibuja con los contornos de las hojas y un color proporcional a $\ell_L$,
como en la Figura 3.1. Da una vista global en la que los patrones locales aparecen
como rectángulos de color distinto dentro de bloques que, en el tablero de ST-CC, se
veían uniformes.

\subsection{4. Mapas de profundidad}

Para cada región se dibuja la profundidad máxima que alcanza el refinamiento en sus
filas, y análogamente para cada intervalo de tiempo. Si la región 5 aparece en hojas
de profundidades 1, 3 y 2, su valor en el mapa es 3. Un mapa con profundidad uniforme
indicaría que el tablero de ST-CC era suficiente. Las zonas con profundidades altas
señalan las regiones y los horarios donde se concentra la interacción
espacio-temporal local, es decir, dónde y cuándo la partición base es insuficiente.

\subsection{5. Frecuencia de soporte}

BBAC$_I$ converge a mínimos locales y el refinamiento es voraz, así que el resultado
puede cambiar con la inicialización. A cada hoja se le asigna su frecuencia de
soporte, que es la proporción de ejecuciones, con distintas inicializaciones de la
partición base, en que aparece una hoja con índice de Jaccard sobre celdas
$\ge\kappa$ respecto de ella.

El índice de Jaccard entre dos conjuntos de celdas es $|A\cap B|/|A\cup B|$. Vale 1
si son idénticos y 0 si no comparten nada.

\paragraph{Ejemplo.} Una hoja de $3\times2$ regiones y horarios tiene 6 celdas. En
otra ejecución aparece una hoja de $2\times2=4$ celdas contenida en la primera. Entonces
$|A\cap B|=4$, $|A\cup B|=6$ y $J=4/6\approx0.67$. Con $\kappa=0.7$ no se cuenta como la
misma hoja, y con $\kappa=0.5$ sí. Si la hoja aparece en 170 de 200 ejecuciones, su
soporte es $0.85$. Una hoja con soporte bajo no debe interpretarse como hallazgo.

\paragraph{Estado de la implementación.} El cálculo de la frecuencia de soporte y del
índice de Jaccard no aparece todavía en \texttt{Implementaciones}. Si el capítulo de
resultados los va a reportar, hay que implementarlos. El valor $\kappa=0.5$ y el
criterio de soporte bajo (menos de la mitad de las ejecuciones) son propuestas que
quedan por confirmar. Además, el comentario de \texttt{AST-CC/02\_refinamiento.py}
asigna la ecuación del nivel de actividad a la versión respecto del padre, cuando en la
tesis esa ecuación es la global.

\begin{idea}
Al abandonar el tablero se pierde la lectura sencilla de ST-CC, y los cinco productos
la reconstruyen. Los descriptores dicen qué es cada hoja, el árbol anotado cómo se
llegó a ella, la matriz teselada dónde está, los mapas de profundidad dónde y cuándo
el tablero fue insuficiente, y la frecuencia de soporte cuánto se puede confiar en
que no sea un accidente de la inicialización.
\end{idea}

% =====================================================================
\part{Capítulo 3, bloque 7: limitaciones y cierre}

\section{Por qué existe esta sección}

Un método propuesto que solo muestra sus virtudes no es defendible. La sección 3.10
declara antes de los experimentos qué no puede hacer el diseño, cómo condiciona la
lectura de los resultados y cómo se piensa medir el efecto de cada limitación. Por eso
cada limitación viene con su forma de evaluarla. La sección termina con las
extensiones que se consideraron y quedan fuera del alcance.

\section{3.10 Las ocho limitaciones}

\paragraph{1. Restricción guillotina.} El refinamiento solo mira hacia dentro de cada
bloque, así que un patrón que cruce la frontera entre dos bloques base no se recupera
como una hoja sino como hojas separadas. Por ejemplo, un patrón en las regiones 3 a 6 en
las mismas horas, con las regiones 3 y 4 en el bloque base A y las 5 y 6 en el B, se
recupera como dos hojas pequeñas y no como un rectángulo de cuatro regiones. Por eso se
menciona como extensión un paso de fusión de hojas vecinas.

\paragraph{2. Dependencia de la partición base.} El refinamiento parte de un óptimo
local de BBAC$_I$, que tiene muchos mínimos locales. La búsqueda local, las múltiples
inicializaciones y la frecuencia de soporte atenúan y miden el efecto. La base también
importa, pues con C5 y con C2 se obtienen tableros muy distintos, y C5 deja menos
interacción local por refinar, porque agrupa las regiones por la forma de su perfil y no
por su volumen.

\paragraph{3. Resolución de entrada.} Los cortes agrupan regiones e intervalos pero no
los subdividen, así que el método no puede bajar de la resolución de la matriz, que son
las zonas de taxi (de tamaños muy dispares) y los intervalos de dos horas.

\paragraph{4. Parámetros.} El método no elimina $K_R$ y $K_T$ sino que cambia su
papel. Fijan la resolución inicial, y la final la deciden $\alpha$ y $\delta$, dos
parámetros con interpretación directa. Esto solo se sostiene si el resultado es
razonablemente estable frente a $K_R$ y $K_T$, lo que debe comprobarse empíricamente.

\paragraph{5. Comparación con rejillas más finas.} Cabe preguntar por qué no basta con
subir $K_R$ y $K_T$ o elegirlos con Cross-Associations. La respuesta debe ser empírica, y
consiste en comprobar que una rejilla global con el mismo número de parámetros, o con el
mismo BIC, no recupera los patrones locales que encuentra AST-CC.

\emph{Ejemplo del conteo de parámetros.} Sea $n=12$ intervalos repartidos en $K_T=4$
grupos de 3 columnas. Un corte de filas en un único bloque de 3 columnas añade
$|J|-1=2$ parámetros. Partir globalmente el mismo grupo de filas a lo largo de los 4
grupos de columnas añade $4\cdot2=\mathbf{8}$. Los parámetros crecen con la estructura
local encontrada y no con la resolución global.

\paragraph{6. Validez de las pruebas.} La selección de nodos con los mismos datos puede
inflar el error de tipo~I. El \emph{bootstrap} de días supone días independientes y usa
$E^B$ como estimación de las medias. Con pocos días, como en el feriado, es algo liberal
en los bloques nocturnos (hasta 0.30 en el peor bloque de la calibración). Por eso es
indispensable un experimento de calibración sobre datos nulos.

\paragraph{7. Costo frente a resolución.} El reparto del nivel por masa da a los
bloques con pocos viajes niveles menores que la resolución alcanzable con $B_{\max}$
réplicas.

\emph{Ejemplo.} Con $N=10^6$ y $\alpha=0.05$, un bloque de 600 viajes recibe
$\alpha_B=0.05\cdot600/10^6=3\cdot10^{-5}$. Supera $N_{\min}=500$, pero el p-valor
mínimo con 20\,000 réplicas es $1/20\,001\approx5\cdot10^{-5}>3\cdot10^{-5}$, así que
queda como hoja inalcanzable. Eso no es lo mismo que homogénea, porque el método no
afirma que no haya nada sino que no puede probarlo al nivel asignado.

\paragraph{8. Significativo no es interesante.} Una hoja puede tener heterogeneidad
detectable y ser irrelevante en el dominio. Por eso se reporta siempre el tamaño de
efecto junto al p-valor, y los patrones principales se contrastan con conocimiento del
dominio, como el foco alrededor de Penn Station que identifican Liu y colaboradores.

\section{Extensiones que quedan como trabajo futuro}

Cada una se consideró durante el diseño y tiene un motivo para quedar fuera.

\begin{itemize}
\item \textbf{Control de FDR jerárquico (Yekutieli).} Es más permisivo que la FWER
cuando se esperan muchos refinamientos. Se dejó fuera porque la FWER es la garantía más
simple de interpretar en una primera versión del método.
\item \textbf{Elegir y probar con días distintos.} Elegir la partición y los cortes con
la mitad de los días y probarlos con la otra mitad elimina el doble uso de los datos a
costa de potencia, en la línea del \emph{count splitting}. Exige comprobar antes que la
mitad de los días basta para estimar la partición.
\item \textbf{Nulas paramétricas por celda.} Con una función de varianza como la ley de
potencia de Taylor. No se adoptaron porque no reproducen las desviaciones correlacionadas
entre las celdas de un bloque (sección 3.5).
\item \textbf{Cortes de columnas en intervalos consecutivos.} Con doce intervalos hay
$2^{11}-1=2047$ biparticiones, pero solo 11 si deben respetar el orden del día. Serían
más fáciles de interpretar y más baratos, a costa de no capturar horarios no
contiguos, como las dos horas punta del día.
\item \textbf{Fusión de hojas vecinas.} Recuperaría los patrones que cruzan bloques base,
que es la limitación de la restricción guillotina.
\end{itemize}

\section{Consideraciones finales del capítulo}

El capítulo presentó AST-CC, una extensión de ST-CC que refina localmente la partición
en tablero mediante un árbol de cortes guiados por pruebas de hipótesis. La pieza
teórica central es la proposición de descomposición por hojas, según la cual la
divergencia del árbol se reparte exactamente entre las hojas y la ganancia de cada corte
es un estadístico de razón de verosimilitud. El corolario muestra que el método
descompone el término de interacción local que ST-CC deja sin modelar, de modo que el
criterio de refinamiento y la función objetivo de ST-CC son la misma cantidad vista a
distintas escalas.

Sobre esa base, el método combina una prueba ómnibus calibrada con \emph{bootstrap} de
los días observados, pruebas dirigidas con ajuste \emph{min-P}, un reparto jerárquico
del nivel y un umbral de tamaño de efecto corregido por sesgo. Cada limitación tiene en
el Capítulo~4 una forma de evaluarse, que es la calibración sobre datos nulos para la
validez de las pruebas, la comparación con rejillas globales de igual número de
parámetros para el valor del refinamiento local, y el análisis de sensibilidad a $K_R$,
$K_T$, $\alpha$ y $\delta$ para la dependencia de la partición base y de los parámetros.

\begin{idea}
AST-CC conserva el tablero de ST-CC y le añade un refinamiento local cuyas decisiones
son pruebas de hipótesis con error controlado. Sus límites son la restricción a cortes
dentro de cada bloque, la dependencia de la partición base y la validez solo aproximada
de las pruebas cuando el árbol se construye con los datos, y cada uno tiene un
experimento previsto que lo mide.
\end{idea}


\end{document}
